% !TEX program = xelatex
% 电工电子技术 · 第二章第一册（Chapter Two-1）· 精讲讲义（lecture-deepdive 规范）
\documentclass[11pt]{ctexart}
%=====================================================================
% theme-violet.tex —— 讲义精讲 · 紫花主题
% 用法：\documentclass[11pt]{ctexart}
%       \input{assets/theme-violet}
%       \begin{document} ... \end{document}
% 只用 ctexart 自带能力 + xcolor/enumitem/fancyhdr/pifont/amsmath，
% 不依赖 tikz/tcolorbox，保证 MiKTeX 可直接编译。
%=====================================================================

\usepackage[a4paper,left=2.0cm,right=2.0cm,top=2.1cm,bottom=2.3cm,%
            headsep=12pt,footskip=22pt]{geometry}
\usepackage{amsmath,amssymb}
\usepackage{xcolor}
\usepackage{array}
\usepackage{enumitem}
\usepackage{pifont}
\usepackage{fancyhdr}
\usepackage[hidelinks]{hyperref}

%---------------------------------------------------------------- 配色
\definecolor{violet}{RGB}{107,78,143}      % 主色：紫
\definecolor{violetdeep}{RGB}{63,44,92}    % 深紫：标题/强调
\definecolor{petal}{RGB}{183,159,208}      % 花瓣：装饰线、花饰
\definecolor{bandviolet}{RGB}{247,243,251} % 浅紫底：档案/提示框
\definecolor{bandcream}{RGB}{252,249,242}  % 米白底：易错/注意框
\definecolor{rulepurple}{RGB}{217,207,230} % 分隔线
\definecolor{inkgray}{RGB}{92,96,110}      % 次要文字

%---------------------------------------------------------------- 字体
% 有 Noto CJK 用 Noto，没有就退回 ctex 默认（Windows 下为宋体/黑体）
\IfFontExistsTF{Noto Serif CJK SC}{\setCJKmainfont{Noto Serif CJK SC}}{}
\IfFontExistsTF{Noto Sans CJK SC}{\setCJKsansfont{Noto Sans CJK SC}}{}

%---------------------------------------------------------------- 版式
\setlist{nosep,leftmargin=1.9em}
\linespread{1.28}
\setlength{\parskip}{3.5pt}
\setlength{\parindent}{0pt}
\setlength{\tabcolsep}{5pt}
\renewcommand{\arraystretch}{1.28}

%---------------------------------------------------------------- 花饰
% \ding{95} = 五瓣小花，\ding{96} = 星形花；都来自 pifont，必定存在
\newcommand{\orn}{\textcolor{petal}{\ding{95}}}          % 行内小标
\newcommand{\ornstar}{\textcolor{petal}{\ding{96}}}      % 星形花
\newcommand{\flowerrule}{%                                % 居中花饰分隔线
  \par\vspace{7pt}%
  \noindent\makebox[\linewidth]{\color{petal}%
    \rule{0.20\linewidth}{0.6pt}\hspace{8pt}{\small\ding{95}}\hspace{4pt}%
    \raisebox{1.1pt}{\tiny\ding{95}}\hspace{4pt}{\small\ding{95}}\hspace{8pt}%
    \rule{0.20\linewidth}{0.6pt}}%
  \par\vspace{7pt}}

%---------------------------------------------------------------- 标题与节头
\newcommand{\doctitle}[3]{%   #1 主标题  #2 副标题  #3 元信息（可用 \\ 分行）
  \begin{center}
    {\color{petal}\small \orn\enspace #2\enspace \ornstar}\\[7pt]
    {\LARGE\bfseries\color{violetdeep} #1}\\[10pt]
    {\color{petal}\rule{0.55\linewidth}{1pt}}\\[4pt]
    {\small\color{inkgray} #3}
  \end{center}\vspace{2pt}}

\newcommand{\hdr}[2]{%   #1 章节号（如 §0）  #2 标题
  \par\vspace{12pt}%
  \noindent\colorbox{violetdeep}{\color{white}\bfseries\small\;#1\;}\hspace{7pt}%
  {\color{violetdeep}\large\bfseries #2}\par\nobreak
  \vspace{3pt}{\color{petal}\hrule height 1.2pt}\vspace{6pt}}

\newcommand{\subhdr}[1]{%
  \par\vspace{6pt}%
  \noindent{\color{violet}\bfseries #1}\par\vspace{1pt}}

%---------------------------------------------------------------- 提示框
\newcommand{\callout}[1]{% 浅紫底：讲义档案、读法说明
  \par\medskip\noindent
  \fcolorbox{rulepurple}{bandviolet}{\parbox{\dimexpr\linewidth-2\fboxsep-2\fboxrule}{\vspace{2pt}#1\vspace{2pt}}}%
  \par\medskip}

\newcommand{\warnbox}[1]{% 米白底：易错、边界条件
  \par\medskip\noindent
  \fcolorbox{rulepurple}{bandcream}{\parbox{\dimexpr\linewidth-2\fboxsep-2\fboxrule}{\vspace{2pt}#1\vspace{2pt}}}%
  \par\medskip}

\newcommand{\gapitem}[2]{% #1 缺口类型（紫） #2 内容
  \par\vspace{2pt}\noindent
  \hangindent=1.4em\hangafter=1
  {\color{petal}\small\ding{96}}\enspace{\bfseries\color{violetdeep}#1}\enspace #2\par}

% 前提旁注：只放在「需要理解的那一步」后面，给陈述 + 本题代入，不重证。
\newcommand{\aside}[1]{%
  \par\vspace{1pt}\noindent
  \hangindent=1.6em\hangafter=1
  {\small\color{inkgray}%
    \hspace*{0.6em}{\color{petal}$\triangleright$}\hspace{2pt}#1}%
  \par\vspace{2pt}}

%---------------------------------------------------------------- 表格
% 本机 MiKTeX 的 colortbl 与 array 版本不匹配（\rowcolor 会报错），
% 故表格统一用「上粗紫线 + 表头 + 中细线」的三线式，不依赖 colortbl。
\newcommand{\tabtop}{\par\vspace{3pt}{\color{violet}\hrule height 1.2pt}\vspace{1pt}}
\newcommand{\tabmid}{\par\vspace{1pt}{\color{petal}\hrule height 0.5pt}\vspace{1pt}}
\newcommand{\tabbot}{\par\vspace{3pt}{\color{rulepurple}\hrule height 0.5pt}\vspace{8pt}}
\newcommand{\thead}[1]{{\bfseries\color{violetdeep}#1}}

%---------------------------------------------------------------- 页眉页脚
\pagestyle{fancy}
\fancyhf{}
\renewcommand{\headrulewidth}{0.4pt}
\renewcommand{\headrule}{\hbox to\headwidth{\color{rulepurple}\leaders\hrule height \headrulewidth\hfill}}
\fancyhead[L]{\small\color{inkgray}\songti \doctag}
\fancyhead[R]{\small\color{inkgray}\songti \coursetag}
\fancyfoot[C]{\small\color{inkgray}\orn\enspace{\color{violet}\thepage}\enspace\orn}

% 每篇讲义在导言区设这两个值
\newcommand{\doctag}{}
\newcommand{\coursetag}{}
\newcommand{\settags}[2]{\renewcommand{\doctag}{#1}\renewcommand{\coursetag}{#2}}

%---------------------------------------------------------------- 数学记号
\newcommand{\vect}[1]{\boldsymbol{#1}}
\newcommand{\dd}{\,\mathrm{d}}
\newcommand{\ddx}{\frac{\mathrm{d}}{\mathrm{d}x}}
\addtolength{\jot}{2pt}



\settags{C2-1 · 精讲}{电工电子技术}

\begin{document}

\doctitle{2.1--2.4 正弦交流电路的基本概念与相量法}
{有效值为什么除以 $\sqrt{2}$、相量凭什么代替正弦量、相量加减与瞬时值加减差在哪}
{依据：\texttt{Chapter Two-1-2026中白上传.pdf}（p1--p43，共 43 页；本讲义页码一律指本册 PDF 物理页，引用他册时写成 [C2-2 p3] 的形式）\ \textperiodcentered\ 2026-10-23}

%=============================================================
\hdr{§0}{课件坐标、记号与三处拼接}

\subhdr{逐页核对的结果：本册实际讲了四段，不是两段}

我按页眉逐页核过 43 页：粗测地图说「p13--p43 全是 2.2」，这个说法不对。
\textbf{2.2.2 在 p27 就结束了}，p28 的页眉已经换成「2.3 Single parameter AC circuit」，
p34 又换成「2.4 Series AC circuits」。所以本册是四段：基本概念、相量表示法、单元件伏安关系、串联 RLC 与阻抗。
四段共用同一套记号与同一个工具（复数），因此这一册可以一次讲完。

\par\vspace{2pt}
\noindent\begin{tabular}{@{}p{5.5cm}p{1.9cm}p{8.1cm}@{}}
\tabtop
\thead{课件小节（页眉原文）} & \thead{页码} & \thead{本讲义对应}
\\
\tabmid
2.1.1 Introduction to AC circuit & p2--p5 & §1.0（「哪一个不是交流」+ 发电机模型）
\\
2.1.2 Expression of AC & p6 & §2.1（三要素）
\\
2.1.3（1）Rate of change & p7--p8 & §2.2（$T$、$f$、$\omega$ 与 314 的来历）
\\
2.1.3（2）Magnitude & p9 & §1（有效值与最大值的关系）
\\
2.1.3（3）Phase and initial phase & p10--p12 & §2.3（E1/E2/E3 三题全解）
\\
2.2.1 Sine wave 与旋转矢量、相量 & p13--p19 & §3（相量为什么能代替正弦量）
\\
2.2.1 复数四形式与四则运算 & p20--p23 & §4（E1/E2/E3 三题全解）
\\
2.2.2 Phasor representation of sine AC & p24--p27 & §5（Example3 全程，主例题）
\\
2.3 Single parameter AC circuit & p28--p33 & §6（R、L、C 三条伏安关系的推导）
\\
2.4.1 Series RLC circuits & p34--p43 & §7（相量 KVL、阻抗、阻抗三角形）
\tabbot
\end{tabular}

\callout{\textbf{三处与粗测地图不同，按版面改过来}：\\
（1）\textbf{2.3 从 p28 起、2.4 从 p34 起}，不是「p13--p43 全是 2.2」。判据是页眉：p27 页眉仍是
「2.2.2 Phasor representation of sinusoidal AC」，p28 页眉已换成「2.3 Single parameter AC circuit」，
p34 换成「2.4 Series AC circuits / 2.4.1 Series RLC circuits」。\\
（2）\textbf{p1、p13、p24、p28、p34 是五张封面页}，只有标题、校徽与
「Electrical Technology DALIAN UNIVERSITY OF TECHNOLOGY」，不承载任何内容；
全册实际承载内容的只有 38 页。\\
（3）\textbf{p24--p43 的 2.3、2.4 在 C2-2 里会以更完整的形式重讲}（C2-2 的 p1--p6 就是 2.4.1 的重复），
本讲义只讲到「用本册的话能自洽」为止：够你读懂 p30--p43 每一条式子怎么来，不替 C2-2 讲串并联与功率。
另：C2-2 §1.1 把「同频率」这条前提引作 [C2-1 p24]，核对后 p24 是 2.2.2 的封面页，
该前提实在本册 \textbf{p14}（§3.3），元件伏安关系实在 \textbf{p29--p33}；本讲义页码一律以页眉为准。}

\callout{\textbf{骨架为什么不照技能模板的「方法 A / 方法 B」排}：课件本身是
「概念（p2--p12）→ 工具（相量 p13--p27、复数 p20--p23）→ 元件（p28--p33）→ 电路（p34--p43）」的四级递进。
相量法与复数法不是并列的两条主线，而是同一条链上的两环（相量是复数，复数是工具），
所以 §3 讲「相量凭什么成立」、§4 讲「复数怎么算」，§5 才把两环扣在一起做主例题。
把主例题放在工具之后，就是这条链的位置；换任何别的排法都会出现「用了还没讲的工具」。}

\subhdr{记号：先统一，后文不再回头}

\par\vspace{1pt}{\small
\noindent\begin{tabular}{@{}p{3.0cm}p{12.6cm}@{}}
\tabtop
\thead{记号} & \thead{约定}
\\
\tabmid
$i,\ u,\ e$ & \textbf{瞬时值}（小写）：随 $t$ 变的那个量，如 $i=12.7\sqrt2\sin(314t+30^\circ)\,\mathrm{A}$
\\
$I_m,\ U_m,\ E_m$ & \textbf{幅值}（最大值、峰值）：$\sin$ 前面那个正系数，恒为正
\\
$I,\ U,\ E$ & \textbf{有效值}（大写、不带下标 $m$）：凡写 $U=220\,\mathrm{V}$ 都是有效值，不指幅值
\\
$\dot I,\ \dot U$ & \textbf{相量}：大写加点，含「模 + 初相」两个信息；本讲义一律用有效值相量，$\dot I=I\angle\psi_i$
\\
$\psi_i,\ \psi_u$ & \textbf{初相}：$t=0$ 时的相位，单位度；$(\omega t+\psi)$ 整体叫\textbf{相位}
\\
$\varphi$ & \textbf{相位差} $\varphi=\psi_u-\psi_i$，同时就是\textbf{阻抗角}（p12 与 p41--p43 是同一个量）
\\
$\omega,\ f,\ T$ & 角频率、频率、周期，$\omega=2\pi f=2\pi/T$
\\
$X_L,\ X_C,\ X$ & 感抗 $X_L=\omega L$、容抗 $X_C=1/(\omega C)$；电抗 $X=X_L-X_C$，\textbf{可正可负}
\\
$Z,\ |Z|$ & 复阻抗 $Z=R+jX$（\textbf{不加点}，它不代表正弦量）；阻抗模 $|Z|=\sqrt{R^2+X^2}$
\\
$j$ & 虚数单位。本课程用 $j$ 不用 $i$，因为 $i$ 已经给了电流
\tabbot
\end{tabular}}

\subhdr{三处拼接（课件自己的缝）}

\gapitem{拼接缺口}{p9 只写「$I=I_m/\sqrt2$」，p25 只写「$\dot I=\sqrt2\angle45^\circ$」，
两页都不写「相量里装的到底是有效值还是幅值」。约定不声明，$\sqrt2$ 就会在每一步漏乘或多乘。
本讲义 §3.5 把它钉成一条：$\dot I=I\angle\psi_i\iff i=\sqrt2I\sin(\omega t+\psi_i)$，再用 p26 的 Example3 反验。}

\gapitem{拼接缺口}{p27 讲「旋转因子」$e^{j90^\circ}=j$，p31 给「$\dot U=jX_L\dot I$」，p32 给「$\dot U=-jX_C\dot I$」。
三页用的是\textbf{同一个} $j$——乘 $j$ 就是逆时针转 $90^\circ$，也就是电压超前电流 $90^\circ$。
课件没有任何一页把这条缝缝上，初学者只会看到 $j$ 莫名其妙地出现在电感的式子里。
本讲义 §6.3、§6.4 在推导末尾把它指出来。}

\gapitem{拼接缺口}{p35 出一题并给了一个被打红叉的错误答案 $U=U_R+U_L+U_C=3\,\mathrm{V}$，
p37 用相量图给出正确答案，p42、p43 才讲 $\varphi$ 与阻抗三角形。三段各说一半：
「为什么不能代数相加」在 p35 没答，「答案为什么是 $1\,\mathrm{V}$」在 p37 只靠图没写式，
「$\varphi$ 是什么」到 p42 才出现。本讲义 §5、§7.2、§7.6 把三段接成一条线。}

\callout{\textbf{本讲义怎么用}：\\
第一遍只看 §0 的表、§5 主例题、§8 公式卡——这三处能让你把题做对；\\
第二遍再看 §1、§3、§6——这三处解释「$\sqrt2$ 从哪来」「相量凭什么能代替正弦量」「电感的 $90^\circ$ 从哪来」，
它们决定你换一个电路还能不能算对；\\
§10 缺口清单逐条对应课件某页某句，课件哪句话让你卡住，就去那里找。}

%=============================================================
\hdr{§1}{有效值：为什么是幅值除以 $\sqrt2$}

\subhdr{1.0 先答 p2 的提问：哪一个不是交流}

课件 p2 给了五个并列波形并问「Which one is not AC?」，页面上只在图组下方写了「pulsing DC」，
\textbf{没有指出是哪一幅、也没有给判据}。判据其实只有一条，就藏在 p2 自己的定义句里：

\par\vspace{2pt}
\noindent\begin{tabular}{@{}p{1.0cm}p{6.4cm}p{7.9cm}@{}}
\tabtop
\thead{图} & \thead{波形（读图记录）} & \thead{是不是交流}
\\
\tabmid
1 & 三角波，正峰、负谷都出现 & 是（跨横轴改变极性）
\\
2 & 单向矩形脉冲串，只在横轴上方，脉冲间回到零 & \textbf{不是}：极性从不改变，是脉动直流
\\
3 & 矩形方波，正负交替 & 是（跨横轴）
\\
4 & 正弦波，正负对称 & 是（本册的主角）
\\
5 & 单向半波脉动，只有一串拱形、无负值 & \textbf{不是}：脉动直流
\tabbot
\end{tabular}

按 p2 的定义句「voltage or current that changes polarity or direction, respectively, over time」，
判据就是\textbf{波形有没有跨过横轴}。2 与 5 都没有跨过，故都不是交流。
\aside{[待确认] p2 图组下方只写了一个「pulsing DC」标签，而 2 与 5 两幅都满足这个描述；
标签究竟指哪一幅、课上是否只认一幅，文字层与读图记录都不能判定。本讲义按判据给出两幅，不作单选。}
另需注意定义句里那个「respectively」：\textbf{电压改的是极性（polarity），电流改的是方向（direction）}，
这是两个说法，不是一回事——译成中文时经常丢掉这层对应。

\subhdr{1.0 后半：p3--p5 的发电机模型，线圈转一圈为什么出来正弦}

p3 的图上标了六个术语（magnet / rotating winding / cut magnetic lines / EMF / brush and ring /
dynamic contact），p4 是同一模型的放大图（去掉全部文字），p5 在同一个模型上补了外电路的电流方向。
三页连起来是一条完整链条，逐项对上：

\par\vspace{2pt}
\noindent\begin{tabular}{@{}p{4.6cm}p{11.0cm}@{}}
\tabtop
\thead{课件术语} & \thead{它在链条里干的事}
\\
\tabmid
magnet（磁铁） & 提供恒定磁场，极间由 N 指向 S
\\
rotating winding（绕组/线圈） & 在磁场中匀速旋转，是「动」的那一方
\\
cut magnetic lines（切割磁力线） & 线圈边与磁力线的夹角随时间变，所以切割速率也随时间变
\\
EMF（电动势） & 切割速率决定感应电动势的大小，方向由右手定则定
\\
brush and ring（碳刷滑环） & 动态接触：把\textbf{旋转}的线圈两端接到\textbf{静止}的外电路
\\
外电路 & 电流经碳刷流出、经仪表（负载）再流回，构成闭合回路
\tabbot
\end{tabular}

\textbf{为什么出来的正好是正弦}（p3--p5 只说「大小与方向一直在变」，没给这一步）：
设线圈平面与磁场的夹角为 $\theta=\omega t$，则穿过线圈的磁通

\[ \Phi=\Phi_m\cos\omega t ,\qquad
   e=-\frac{\dd\Phi}{\dd t}=\omega\Phi_m\sin\omega t . \]

求导一次，$\cos$ 变 $\sin$，并且\textbf{多出一个 $\omega$}。所以「大小与方向一直在变」不是形容，
而是「匀速旋转 + 余弦磁通」的必然结果。这也正是 p6 那条 $i=I_m\sin(\omega t+\psi)$ 的物理出生地，
以及 §3.4 那条 $\dd/\dd t\leftrightarrow j\omega$ 的物理含义：求导会多出一个 $\omega$、并把相位推 $90^\circ$。

线圈每转半圈，切割方向反一次，电动势的符号就反一次——这就是交流。
p3 右下角「switching polarity 改变极性 / switching direction 改变方向」说的正是这两件事；
p5 的对比句「The current only flows in one direction in DC circuit. But in AC the current changes
direction periodically」则是把它与直流分开。

\aside{真实发电机的电角频率等于极对数乘机械角速度；本册只讨论电角频率 $\omega$，
不再区分极对数——凡出现 $\omega$ 一律按电角频率读。}

\subhdr{1.1 三个「大小」的分工（p9 那张表）}

p9 左边列了三行符号，右边各贴一个名字。三者的区别不在符号长短，而在「什么时候用」：

\par\vspace{2pt}
\noindent\begin{tabular}{@{}p{2.4cm}p{3.6cm}p{9.6cm}@{}}
\tabtop
\thead{记号} & \thead{名字} & \thead{什么时候用}
\\
\tabmid
$e,\ i,\ u$ & 瞬时值 & 写波形、代入某个具体时刻、做瞬时值相加（§5.4）
\\
$E_m,\ I_m,\ U_m$ & 幅值（最大值、峰值） & 只出现在 $\sin$ 前面；示波器看峰值得用它
\\
$E,\ I,\ U$ & 有效值 & \textbf{一切仪表读数、铭牌数据、相量的模}；$220\,\mathrm{V}$ 指的就是它
\tabbot
\end{tabular}

一句话记住：\textbf{有效值是「能用直流表量出来的那个数」，幅值是「波形顶到多高」}。
两者差一个 $\sqrt2$，而这个 $\sqrt2$ 不是规定，是算出来的。

\subhdr{1.2 有效值的定义：热等效，以及那一步积分}

p9 的蓝框标题是「Relationship between effective value and amplitude」，框里直接写
$I=I_m/\sqrt2$、$U=U_m/\sqrt2$、$E=E_m/\sqrt2$，\textbf{三个结果一个推导都没有}。
这正是本册第一个必须讲透的机制。补出来只用两步。

\textbf{第一步：定义。}有效值的定义不是「平均」，而是\textbf{热等效}：让直流 $I$ 与交流 $i$
分别通过同一个电阻 $R$，在一个周期 $T$ 内产生的热量相等，就称这个直流值 $I$ 为该交流的有效值。

\[ I^2RT=\int_0^T i^2R\dd t
   \quad\Longrightarrow\quad
   I=\sqrt{\frac{1}{T}\int_0^T i^2\dd t} . \]

所以有效值的学名叫\textbf{均方根值}（root mean square）：先平方、再取一个周期的平均、最后开方。

\textbf{第二步：算。}取 $i=I_m\sin\omega t$（先令 $\psi=0$，初相不影响这一定义的取值）：

\[ i^2=I_m^2\sin^2\omega t=\frac{I_m^2}{2}\left(1-\cos2\omega t\right) . \]

一个周期内 $\cos2\omega t$ 的平均值为零（因为 $\omega T=2\pi$，正好走完两个整周期）：

\[ \frac{1}{T}\int_0^T\cos2\omega t\dd t=0 ,\qquad
   \frac{1}{T}\int_0^T i^2\dd t=\frac{I_m^2}{2} , \]

于是

\[ I=\sqrt{\frac{I_m^2}{2}}=\frac{I_m}{\sqrt2}\approx0.707\,I_m . \]

反过来就是 p9 那三行的用法：

\[ I_m=\sqrt2\,I\approx1.414\,I . \]

\aside{p9 只印结果、不印这一步，这就是 §10 清单的第一条（机制缺口）。
「$0.707$」这个系数只对正弦波成立；方波、三角波的有效值系数不同，本册不做，遇到要另算。}

\warnbox{\textbf{为什么不能取平均值}：正弦量在一个整周期上的平均值为零，
\[ \frac{1}{T}\int_0^T I_m\sin\omega t\dd t=0 , \]
它连「大小」都不反映。所以不能拿平均值当有效值——这就是「必须先平方」的全部理由。}

\subhdr{1.3 220 V 与 311 V 是同一个电压}

p9 右下角那张照片里，万用表插在插座上，读数 220；照片上方问了两句：
「1. 220V？Effective or amplitude (peak value)?」与
「2. Is the reading of the meters amplitude or effective value?」
\textbf{两问都没有答案}（第二问在 p30 又问了一次）。答案是：

\begin{itemize}
\item 市电 $220\,\mathrm{V}$ 是\textbf{有效值}；它的幅值是
\[ U_m=\sqrt2\times220\approx311\,\mathrm{V} , \]
这正是 p11 的 E2 里那个 $311$ 的出处：$u=311\sin(314t-60^\circ)\,\mathrm{V}$。
所以「220 V」与「311 V」是同一个电压的两种说法，不是两个电压。
\item 万用表（交流电压挡）读的是\textbf{有效值}。它的表头按正弦有效值刻度，
所以拿它去量方波或可控硅调光后的波形会读错——那些波形的有效值系数不是 $\sqrt2$。
\item 有效值不做特殊说明时恒指正弦量。p26 的 Example3 全部用有效值写相量，
p9 那条 $I=I_m/\sqrt2$ 就是它每一次转换的依据。
\end{itemize}

\aside{后文凡写 $U$、$I$（不带下标 $m$）一律是有效值；$\dot U$、$\dot I$ 的模也是有效值。
代入本题的形式：$u=311\sin(314t-60^\circ)\,\mathrm{V}$ 的相量是 $\dot U=220\angle-60^\circ\,\mathrm{V}$，
不是 $311\angle-60^\circ$。}

\subhdr{1.4 有效值不能直接相加（这一条预告 §5）}

p35 的提问是「$U_R=1\,\mathrm{V}$、$U_L=1\,\mathrm{V}$、$U_C=1\,\mathrm{V}$，求总电压」，
课件给出的 $U=1+1+1=3\,\mathrm{V}$ 被画了一个红叉。红叉打对了，
但\textbf{为什么错、正确的数是多少，p35 一页没说}。正确的值是 $1\,\mathrm{V}$，理由在 §7.2。

先在这里立一条判据：\textbf{有效值相加只在「同相位」时才等于代数相加}。
p26 的 Example3 有两个同频电流，有效值 $12.7\,\mathrm{A}$ 与 $11\,\mathrm{A}$，
但结果是 $16.8\,\mathrm{A}$ 而不是 $12.7+11=23.7\,\mathrm{A}$——因为两者相位差恰好 $90^\circ$。
这是本册第二个必须讲透的机制，§5 全程演算。

\aside{同一个陷阱的两个面孔：$12.7+11=23.7$（p26 的电流）与 $1+1+1=3$（p35 的电压）。
凡看见「有效值相加」，先问一句「它们同相吗」。}

%=============================================================
\hdr{§2}{三要素与相位差}

\subhdr{2.1 表达式：$i=I_m\sin(\omega t+\psi)$ 里的三个数（p6）}

p6 把正弦量的写法拆成四格：表达式、波形、含义、三要素，最后归纳成
\textbf{幅值、角频率、相位}三要素。逐项对上：

\par\vspace{2pt}
\noindent\begin{tabular}{@{}p{2.8cm}p{2.6cm}p{10.3cm}@{}}
\tabtop
\thead{参数} & \thead{名字} & \thead{含义与其在式中的位置}
\\
\tabmid
$I_m$ & 幅值 & $\sin$ 的系数，恒正；波形顶到多高
\\
$\omega$ & 角频率 & $t$ 的系数，单位 $\mathrm{rad/s}$；波形走多快
\\
$\psi$ & 初相 & 常数项，单位度；$t=0$ 时刻走到哪了
\\
$\omega t+\psi$ & 相位 & 整个括号叫相位，随 $t$ 匀速增长
\tabbot
\end{tabular}

\textbf{一处表述要说清}：p6 把三要素写成「Amplitude / Angular frequency / Phase」，
但右侧含义栏里又写「$\psi$: Initial Phase」。严格说：
\textbf{相位是 $\omega t+\psi$（随时间变），初相是 $\psi$（定值）}。
三要素里第三要素准确的名字是\textbf{初相}。本讲义后文一律用「初相 $\psi$」，不再混称。

\warnbox{\textbf{全册以 $\sin$ 为基准}。p6 到 p43 每一条式子都是 $\sin$，没有一处改用 $\cos$。
如果你从别的书搬来 $\cos$ 形式的公式，必须先把相位换算 $90^\circ$ 再用，否则处处差 $90^\circ$。
本讲义不混用 $\cos$。}

\subhdr{2.2 周期、频率、角频率：p7 那张图要自己验一遍（p7--p8）}

p7 给的是同一个东西的三个说法：

\[ \omega=2\pi f=\frac{2\pi}{T}\ (\mathrm{rad/s}) ,\qquad f=\frac{1}{T} . \]

图上的数是：$T=0.5\,\mathrm{s}$、$f=2\,\mathrm{Hz}$、$\omega=4\pi\,\mathrm{rad/s}$，
横轴上 $1\,\mathrm{s}$ 的跨度里正好两个周期。三个数自洽（自己验）：
$f=1/T=1/0.5=2\,\mathrm{Hz}$；$\omega=2\pi f=4\pi\approx12.57\,\mathrm{rad/s}$。三个数都对。

\textbf{本册最常用的一个换算}：中国电网 $f=50\,\mathrm{Hz}$，则

\[ \omega=2\pi\times50\approx314.16\approx314\,\mathrm{rad/s} . \]

p26 的 Example3 里，两个电流的相位都是 $314t$，
这个 $314$ 就是 $50\,\mathrm{Hz}$ 换成角频率的结果，不是某个元件的参数。
若电网是 $60\,\mathrm{Hz}$（如美国），则 $\omega\approx377\,\mathrm{rad/s}$。

\aside{代入本题的形式：Example3 的 $i_1=12.7\sqrt2\sin(314t+30^\circ)$ 中 $\omega=314\,\mathrm{rad/s}$，
所以 $f=314/(2\pi)\approx50\,\mathrm{Hz}$、$T=1/f\approx0.02\,\mathrm{s}$。
两个电流的 $\omega$ 相同，这是它们能画在同一张相量图上的\textbf{唯一}理由（§3.3）。}

\textbf{p8 的提问「Why 50Hz/60Hz?」没有答案}。下面这段是背景，不属于课件原文：
频率是工程折中的产物。频率高，变压器与电机可以做得小、同样的铁芯能传更多功率，
但铁损、集肤效应与线路感抗压降都随频率上升；频率低，变压器体积大、白炽灯闪烁明显。
19 世纪末两套工业体系分别定了标准：欧洲（西门子、AEG 一侧）取 $50\,\mathrm{Hz}$，
美国（西屋、特斯拉一侧）取 $60\,\mathrm{Hz}$；中国电网沿用 $50\,\mathrm{Hz}$。
这也解释了为什么本册例题里 $\omega$ 都是 $314$ 而不是 $377$。
\textbf{[教材补充]}之所以补这一段：p8 只提问不给答案，而「为什么是 50 Hz」直接决定了所有例题里的 314，
不补的话 314 就成了一个背下来的常数。

\subhdr{2.3 相位、初相与相位差：p10--p12 三题全解}

p10 的定义句是「Phase: the divisional part of a cycle through which the quantity moves forward
from a selected origin」，译过来就是「从选定起点算起，这个量已经走过的、以周期为单位的份额」。
p10--p12 各给一题：

\par\vspace{2pt}
\noindent\begin{tabular}{@{}p{1.1cm}p{5.2cm}p{4.0cm}p{5.4cm}@{}}
\tabtop
\thead{页} & \thead{题} & \thead{相位} & \thead{初相}
\\
\tabmid
p10 & E1：$i=10\sin(1000t+30^\circ)\,\mathrm{A}$ & $1000t+30^\circ$ & $\psi_i=30^\circ$
\\
p11 & E2：$u=311\sin(314t-60^\circ)\,\mathrm{V}$ & $314t-60^\circ$ & $\psi_u=-60^\circ$
\\
p12 & E3：上两式放在一起，求相位差 & —— & $\varphi=\psi_u-\psi_i$
\tabbot
\end{tabular}

p12 的完整演算：

\[ \varphi=\psi_u-\psi_i=-60^\circ-30^\circ=-90^\circ . \]

\textbf{$\varphi<0$ 的含义是「$u$ 滞后 $i$ $90^\circ$」}，等价的说法是「$i$ 超前 $u$ $90^\circ$」。
判读只有一张小表，背下来比记文字快：

\par\vspace{2pt}
\noindent\begin{tabular}{@{}p{3.4cm}p{5.4cm}p{6.9cm}@{}}
\tabtop
\thead{$\varphi=\psi_u-\psi_i$} & \thead{说法} & \thead{波形上的样子}
\\
\tabmid
$\varphi>0$ & $u$ 超前 $i$（$i$ 滞后 $u$） & $u$ 的过零点、峰值都出现在 $i$ 之前
\\
$\varphi<0$ & $u$ 滞后 $i$（$i$ 超前 $u$） & $u$ 的一切都比 $i$ 晚到
\\
$\varphi=0$ & 同相 & 两条曲线同时过零、同时到峰（p30 的电阻）
\\
$|\varphi|=180^\circ$ & 反相 & 一条到正峰时另一条正在负峰
\tabbot
\end{tabular}

\warnbox{\textbf{三条限制，课件没有明说}：
（1）相位差\textbf{只对同频率的两个量有定义}——频率不同，「差」随时间变，没有定值；
（2）初相一般取 $|\psi|\le180^\circ$，超出的部分折回来（如 $\psi=270^\circ$ 写作 $\psi=-90^\circ$）；
（3）$\varphi$ 的正负取决于你把谁减谁，p12 用的是 $\psi_u-\psi_i$，全册不再改，中途也不要换顺序。}

\subhdr{2.4 相位差与阻抗角是同一个量}

p41--p43 上写的是 $\varphi=\psi_u-\psi_i$，与 p12 的相位差完全一样。
这不是巧合：\textbf{「电压比电流超前多少」这件事，在元件层面叫相位差，在整条电路层面就叫阻抗角}。
p31 说「电感上电压超前电流 $90^\circ$」，p42 说「感性电路 $0<\varphi<90^\circ$」，
说的正是同一件事的两个尺度：前者是单个元件（$90^\circ$ 整），后者是 RLC 串起来的整体（$90^\circ$ 以内）。
§7.6 会把这条对应关系写成表。

\aside{记号提醒：课件 p12 的相位差印成 $\phi$，p41--p43 的阻抗角印成 $\varphi$，
文字层里两者转写成同一个字符。本讲义统一写 $\varphi$，不再区分——因为它们是同一个量。}

%=============================================================
\hdr{§3}{相量：一个复数凭什么能代替一条正弦曲线}

\subhdr{3.1 p15 那句话在说什么：旋转矢量的纵轴投影}

p15 的原句是「The vector rotates in an anticlockwise direction if we take the projection of this
vector in Y-axis then we get the instantaneous value of this sinusoidal signal」。
把它写成式子就三件事：长度取幅值、起始角取初相、角速度取 $\omega$。

\begin{itemize}
\item 矢量长度 $=U_m$（幅值）；
\item $t=0$ 时它与横轴夹角 $=\psi$（初相）；
\item 以角速度 $\omega$ 逆时针匀速旋转，故 $t$ 时刻夹角 $=\omega t+\psi$；
\item 取它在\textbf{纵轴}上的投影，就是瞬时值 $u(t)=U_m\sin(\omega t+\psi)$。
\end{itemize}

用复数写，这一整段就是一句：先转成复数、取虚部，

\[ u(t)=\operatorname{Im}\left[U_m e^{j(\omega t+\psi)}\right]=U_m\sin(\omega t+\psi) . \]

\aside{为什么全册都用 $\sin$ 而不用 $\cos$：p15 明确说取\textbf{纵轴}投影。
取纵轴得到 $\sin$，取横轴就得到 $\cos$。p6 到 p43 一律 $\sin$，与这里取的轴一致。
这条看似琐碎，但它决定了 §6 里「电感电压超前 $90^\circ$」还是「滞后 $90^\circ$」——基准轴一换，正负全反。}

\subhdr{3.2 冻结在 $t=0$：从旋转矢量到相量（p17--p19）}

p17 的副标题是「Sine wave $\leftrightarrow$ Rotating vector $\leftrightarrow$ stationary vector」，
p18 又加了一个「phasor」，p19 再加「Complex number」。四步在讲同一件事，只是把时间因子一层层剥掉。

关键一步是把上面那个指数拆开：

\[ U_m e^{j(\omega t+\psi)}=\underbrace{\left(U_m e^{j\psi}\right)}_{\text{不含 }t}\cdot e^{j\omega t} . \]

括号里的复数\textbf{不随时间变}——它就是「冻结在 $t=0$ 的矢量」，也就是 p17 说的 stationary vector。
把这个数除以 $\sqrt2$ 换成有效值、并命名为相量：

\[ \dot U=U e^{j\psi}=U\angle\psi . \]

所以相量是\textbf{一个复数}，不是矢量、不是函数。p18 给的定义句说得很克制：
「A phasor is a mathematical quantity created in electronics to explain AC behavior」——
它是人为了方便解释交流行为而造出来的数学量。

\textbf{被剥掉的是什么？}是 $e^{j\omega t}$，也就是\textbf{频率信息}。
这条很重要：\textbf{相量里查不到频率}。所以
（1）两个相量能不能相加，只能由你在心里记住的 $\omega$ 判断（§3.3）；
（2）由相量还原瞬时值时，必须把 $\omega$ 补回去（§5.2 第 7 步）；
（3）同一张相量图上不允许出现两种频率。

\subhdr{3.3 同频是相量法的门票（p14 的「same frequency」）}

p14 的提问句是「When we have two sinusoidal AC current with the \textbf{same frequency}, different
phase and different amplitude, how to find the total current?」。
那句话里藏着前提，却没说为什么必须有这个前提。证明只有三行：

设 $i_1=\sqrt2I_1\sin(\omega t+\psi_1)$、$i_2=\sqrt2I_2\sin(\omega t+\psi_2)$，相量记为
$\dot I_1=I_1\angle\psi_1$、$\dot I_2=I_2\angle\psi_2$。按 §3.1 的写法，

\[ i_1=\operatorname{Im}\left[\sqrt2\dot I_1e^{j\omega t}\right],\qquad
   i_2=\operatorname{Im}\left[\sqrt2\dot I_2e^{j\omega t}\right] . \]

$\operatorname{Im}$ 是线性运算（取虚部可以逐项进行、可以与常数交换），于是

\[ i_1+i_2=\operatorname{Im}\left[\sqrt2\left(\dot I_1+\dot I_2\right)e^{j\omega t}\right] . \]

右边的相量部分是 $\dot I_1+\dot I_2$，角频率仍是同一个 $\omega$。
\textbf{这一步的唯一支点是「两个 $e^{j\omega t}$ 是同一个东西，才能提出公因子」。}

\textbf{频率不同会怎样}：取 $i_1=\sqrt2\sin314t$、$i_2=\sqrt2\sin1000t$，则
$i_1+i_2=2\sin657t\cos343t$——包络 $2\cos343t$ 随时间慢慢起伏，
整体\textbf{不再是正弦}，没有单一幅值与初相，因此\textbf{没有相量}。
这就是 p26 例题里两个 $314t$ 为什么必须写成一样的理由。

\aside{这张门票的三种失效情形，考试里都出现过：
（1）两个电源频率不同——相量不能相加，只能回时域；
（2）一个直流分量叠在交流上——直流量不属于相量域，要分成两个问题算；
（3）频率随时间变（调频、变频器）——相量法整体失效，要回到瞬时式。}

\subhdr{3.4 微分与积分：全册没写，却是 §6 的唯一工具}

六个元件里有两个的伏安关系是微分形式（$u_L=L\dfrac{\dd i}{\dd t}$、$i_C=C\dfrac{\dd u}{\dd t}$）。
要把它们变成代数式，必须先把「求导」翻译成相量域里的乘法。这一步全册 43 页没有一处写出来，
但 p31、p32 的两条相量式全靠它。规则来自对 $e^{j\omega t}$ 求导：

\[ \frac{\dd}{\dd t}e^{j\omega t}=j\omega\,e^{j\omega t}
   \quad\Longrightarrow\quad
   \frac{\dd}{\dd t}\ \longleftrightarrow\ j\omega ,\qquad
   \int\dd t\ \longleftrightarrow\ \frac{1}{j\omega} . \]

一句话记住：\textbf{相量法的全部力量，来自「微分变成乘一个数」}。
微分方程之所以能变成代数方程，就是因为这条替换。

\aside{代入本题的形式：$u_L=L\dfrac{\dd i}{\dd t}$ 中 $i=\sqrt2I\sin\omega t$，
求导后多出一个 $\omega$、相位多出 $90^\circ$（即乘 $j$），于是 $\dot U_L=j\omega L\dot I$。
$j\omega L$ 就是 $jX_L$——§6.3 会把这步走完。}

\subhdr{3.5 相量里装有效值还是幅值：本册的约定}

这一条不声明，后面每个数都会差 $\sqrt2$。本册的约定（由 p26 反推得出）是\textbf{有效值相量}：

\[ \boxed{\ \dot I=I\angle\psi_i\iff i=\sqrt2I\sin(\omega t+\psi_i)\ } \]

验算一遍就看出为什么：把 $i=\sqrt2I\sin(\omega t+\psi_i)$ 写成取虚部的形式，

\[ i=\sqrt2\operatorname{Im}\left[Ie^{j(\omega t+\psi_i)}\right]
   =\operatorname{Im}\left[\sqrt2\dot I e^{j\omega t}\right] , \]

与 §3.3 的写法完全一致。所以\textbf{相量的模是有效值，还原瞬时值时要在前面补一个 $\sqrt2$}。

p25 的例子正好把这个坑露出来：它写 $\dot I=1+j=\sqrt2\angle45^\circ$，
但没有说对应哪个瞬时式。按本册约定，它对应

\[ i=\sqrt2\times\sqrt2\sin(\omega t+45^\circ)=2\sin(\omega t+45^\circ)\ \mathrm{A} . \]

\textbf{p26 的 Example3 把这个约定钉死了}：题目给 $i_1=12.7\sqrt2\sin(314t+30^\circ)\,\mathrm{A}$，
下一步就写 $\dot I_1=12.7\angle30^\circ\,\mathrm{A}$——$12.7$ 是有效值，$\sqrt2$ 被丢掉了。
如果相量的模取幅值，那就该写 $12.7\sqrt2=17.96\angle30^\circ$，与课件不符。

\warnbox{\textbf{两套标定差 $\sqrt2$ 倍，不能混用}。判据只有一条：
\textbf{看 $\sin$ 前面有没有 $\sqrt2$}。有 $\sqrt2$，相量的模取 $I$；没有 $\sqrt2$（纯幅值写法），
相量的模取 $I_m$。做题时先标一句「本解用有效值相量」，再往下写，能避掉一大半的错。}

%=============================================================
\hdr{§4}{复数工具：四种形式、四则运算、三个例题}

\subhdr{4.1 四种形式与两次改写（p20）}

p20 在一支复平面矢量的旁边竖着写了四个等式，这就是复数的四种形式：

\par\vspace{2pt}
\noindent\begin{tabular}{@{}p{3.1cm}p{5.5cm}p{7.0cm}@{}}
\tabtop
\thead{形式（p20 原名）} & \thead{写法} & \thead{什么时候用它最省事}
\\
\tabmid
直角形式（Rectangular） & $A=a+jb$ & 加减、看实部虚部、代进 $Z=R+jX$
\\
三角形式（Triangular） & $A=c(\cos\psi+j\sin\psi)$ & 从模与角展开成实部虚部
\\
指数形式（Exponential） & $A=c\,e^{j\psi}$ & 乘除、求导、写旋转因子
\\
极坐标形式（Polar） & $A=c\angle\psi$ & 相量图、读相位、写最终答案
\tabbot
\end{tabular}

\textbf{p20 只印结论、没印两次改写。}补出来是两对式子。从直角到极坐标：

\[ c=\sqrt{a^2+b^2} ,\qquad \psi=\arctan\frac{b}{a} ; \]

从极坐标回到直角（三角形式展开即可）：

\[ a=c\cos\psi ,\qquad b=c\sin\psi . \]

\aside{p20 图下的三个注记 real part / imaginary part / modulus 分别对应 $a$、$b$、$c$，
但\textbf{没有一个式子说出 $c$ 与 $a,b$ 的关系}，也没说 $\psi$ 怎么由 $a,b$ 算。
这两对式子就是 §10 清单里的「代数跳步」那条，p21 的 $1+j\leftrightarrow\sqrt2\angle45^\circ$ 全靠它们。}

\warnbox{\textbf{$\arctan$ 只给 $(-90^\circ,90^\circ)$，这个坑必须单独记}。
$a>0$ 时 $\psi=\arctan(b/a)$ 直接可用；$a<0$ 时（点在左半平面）必须加 $180^\circ$。
例：$-1+j$，$\arctan(1/-1)=-45^\circ$ 是错的，正确值是 $-45^\circ+180^\circ=135^\circ$，
即 $-1+j=\sqrt2\angle135^\circ$。验算：$\sqrt2\cos135^\circ=-1$、$\sqrt2\sin135^\circ=+1$，对。}

\subhdr{4.2 p21 的提问：为什么需要这么多种形式}

p21 列了一张对照表（$1+j\leftrightarrow\sqrt2\angle45^\circ$、$j\leftrightarrow1\angle90^\circ$），
底下问「Why do we need so many forms?」，\textbf{没有答案}。答案是一张分工表：

\par\vspace{2pt}
\noindent\begin{tabular}{@{}p{4.2cm}p{4.2cm}p{7.4cm}@{}}
\tabtop
\thead{要干的事} & \thead{用哪种形式} & \thead{理由}
\\
\tabmid
加、减 & 直角形式 & 实部加实部、虚部加虚部，没有别的路
\\
乘、除 & 极坐标（或指数） & 模相乘除、角相加减，一步到位
\\
求相位差、画相量图 & 极坐标 & 相位就是角，直接读
\\
代进 $Z=R+jX$、求实部虚部 & 直角形式 & $R$、$X$ 天然是实部、虚部
\tabbot
\end{tabular}

口诀八个字：\textbf{加减用直角，乘除用极坐标}。§4.3 的四条公式就是这句话的展开。

\subhdr{4.3 四则运算：p22 那四条式子，以及一处括号}

p22 给了四条。前两条没问题，后两条\textbf{课件印的角没有括号}，按字面读会读成别的意思：

\par\vspace{2pt}
\noindent\begin{tabular}{@{}p{2.6cm}p{8.0cm}p{5.2cm}@{}}
\tabtop
\thead{运算} & \thead{正确写法} & \thead{课件 p22 的印法}
\\
\tabmid
加减 & $A_1\pm A_2=(a_1\pm a_2)+j(b_1\pm b_2)$ & 同（无问题）
\\
乘法 & $A_1A_2=c_1c_2\angle(\psi_1+\psi_2)$ & $A_1A_2=c_1c_2\angle\psi_1+\psi_2$
\\
除法 & $\dfrac{A_1}{A_2}=\dfrac{c_1}{c_2}\angle(\psi_1-\psi_2)$ & $\dfrac{A_1}{A_2}=\dfrac{c_1}{c_2}\angle\psi_1-\psi_2$
\tabbot
\end{tabular}

\aside{[待确认] 上表第四行的「课件印法」由\textbf{两路独立读法}得出：文字层抽出的式子与逐页读图记录
都读作 $\angle\psi_1+\psi_2$、角上没有括号。若原页其实有括号而两路都漏读，
则本条退化为「写法提示」，但不影响后面推出的法则本身——法则本身另有推导。}

\textbf{为什么必须加括号}：$\angle$ 的后面是一个角，$\angle(\psi_1+\psi_2)$ 才是「转 $\psi_1+\psi_2$」；
按课件的印法，字面读作「$\angle\psi_1$ 这个复数，再加上一个数 $\psi_2$」，
后者是个无量纲的角，与复数相加根本不合量纲。

\textbf{两条法则的来历}（p22 直接给结果，不推导）：用指数形式相乘，

\[ A_1A_2=\left(c_1e^{j\psi_1}\right)\left(c_2e^{j\psi_2}\right)
   =c_1c_2\,e^{j(\psi_1+\psi_2)} , \]

指数相加就是角相加，模相乘——这就是乘法法则。除法同理（指数相减），
并且要求 $A_2\ne0$，这正是 p22 在 $A_2$ 后面写「$\ne0$」的原因。

\subhdr{4.4 旋转因子：$e^{j90^\circ}=j$ 是什么意思（p27）}

p27 的标题是「rotating factor」，图上把同一个相量 $\dot I$ 转了两个方向：

\[ e^{j90^\circ}=j ,\qquad e^{-j90^\circ}=-j ,\qquad
   \dot I\cdot j=\dot I\,e^{j90^\circ} . \]

按 §4.3 的乘法法则，乘 $j$ 就是「模不变、角加 $90^\circ$」，
所以在复平面上一支相量\textbf{乘 $j$ 就是逆时针转 $90^\circ$，乘 $-j$ 就是顺时针转 $90^\circ$}。
这条为什么重要：§6.3、§6.4 里电感电压写 $\dot U=jX_L\dot I$、电容电压写 $\dot U=-jX_C\dot I$，
那个 $j$ 与 $-j$ 一个不多、一个不少，正好就是「超前 $90^\circ$」与「滞后 $90^\circ$」。
\textbf{同一个符号，在 p27 叫旋转因子，在 p31 叫感抗的一部分}——这就是 §0 列的第二处拼接。

\aside{[待确认] p27 版面上「Question:」后面是空的，第二行右侧那一格只画了坐标轴与虚线、
没有画出 $\dot I\cdot(-j)$ 的箭头。文字层与读图记录都读到空框，故判定为「题目缺失」而非「内容看不清」。
若要补，最自然的填空是：把 $\dot I$ 顺时针转 $90^\circ$，画出指向右下的箭头，标 $\dot I\cdot(-j)=-j\dot I$。}

\subhdr{4.5 p23 三个例题全程（课件只出题、不给答案）}

p23 只有三行题面，\textbf{一个答案都没有}。三题各走两遍（极坐标法与直角法互校），
下面的数已用 Python 复算。

\textbf{E1：$(1+j)+(5+j6)$}。纯加减，用直角形式，实部加实部、虚部加虚部：

\[ (1+j)+(5+j6)=(1+5)+j(1+6)=6+j7 . \]

换成极坐标：$|A|=\sqrt{6^2+7^2}=\sqrt{85}\approx9.22$，$\psi=\arctan(7/6)\approx49.4^\circ$，
即 $6+j7\approx9.22\angle49.4^\circ$。

\textbf{E2：$\dfrac{1+j}{\dfrac{\sqrt3}{2}+j\dfrac12}$}。一眼看不出，先看分母：
$\dfrac{\sqrt3}{2}+j\dfrac12=1\angle30^\circ$（模为 $1$，因为 $\left(\frac{\sqrt3}{2}\right)^2+\left(\frac12\right)^2=1$）。
分子 $1+j=\sqrt2\angle45^\circ$。极坐标法一步到位：

\[ \frac{\sqrt2\angle45^\circ}{1\angle30^\circ}=\sqrt2\angle(45^\circ-30^\circ)=\sqrt2\angle15^\circ . \]

直角法互校（分母乘以共轭，分母模为 $1$，约分最省事）：

\[ \frac{1+j}{\frac{\sqrt3}{2}+j\frac12}
   =\left(1+j\right)\left(\frac{\sqrt3}{2}-j\frac12\right)
   =\left(\frac{\sqrt3}{2}+\frac12\right)+j\left(\frac{\sqrt3}{2}-\frac12\right)
   \approx1.366+j0.366 . \]

两法一致：$\sqrt2\cos15^\circ=1.366$、$\sqrt2\sin15^\circ=0.366$。

\textbf{E3：$\dfrac{(1+j)\left(\dfrac12+j\dfrac{\sqrt3}{2}\right)}{\dfrac{\sqrt3}{2}+j\dfrac12}$}。
分子第二个因子是 $1\angle60^\circ$，分母仍是 $1\angle30^\circ$，分子第一个因子 $\sqrt2\angle45^\circ$。极坐标法：

\[ \frac{\sqrt2\angle45^\circ\times1\angle60^\circ}{1\angle30^\circ}
   =\sqrt2\angle(45^\circ+60^\circ-30^\circ)=\sqrt2\angle75^\circ . \]

直角法互校：

\[ (1+j)\left(\frac12+j\frac{\sqrt3}{2}\right)
   =\left(\frac12-\frac{\sqrt3}{2}\right)+j\left(\frac12+\frac{\sqrt3}{2}\right)
   \approx-0.366+j1.366 , \]

这个中间结果与 $\sqrt2\angle105^\circ$ 相符（$45^\circ+60^\circ=105^\circ$），
再除以 $1\angle30^\circ$ 得 $\sqrt2\angle75^\circ\approx0.366+j1.366$。两法一致。

\warnbox{\textbf{三题的共同教益}：E1 只能直角，E2、E3 极坐标快十倍。
但\textbf{极坐标法只在乘除时成立}——把 E1 拿去做「$9.22\angle49.4^\circ$ 加什么」是算不出来的，
必须退回直角形式。判据就是 §4.2 那张表。}

%=============================================================
\hdr{§5}{主例题：相量加减与瞬时值加减差在哪（p26 Example3）}

\subhdr{5.1 题面与要交出的三个数}

p26 的题面（原文照录）：

\callout{Given that $i_1=12.7\sqrt2\sin(314t+30^\circ)\,\mathrm{A}$，$i_2=11\sqrt2\sin(314t-60^\circ)\,\mathrm{A}$，
Find $i=i_1+i_2$。\\
要交出的三个数是：合成电流的\textbf{有效值} $I$、它的\textbf{初相} $\psi_i$、以及写回\textbf{瞬时式}的结果。}

这三个数分别对应相量的模、角、以及还原后的式子。p26 一路写下去，中间跳了三步：
（1）为什么能把 $12.7\sqrt2$ 直接写成 $12.7$；（2）为什么 $12.7\angle30^\circ$ 与 $11\angle-60^\circ$ 能这样相加；
（3）$16.8\angle-10.9^\circ$ 怎么又变回 $16.8\sqrt2\sin(\cdots)$。下面逐步补齐。

\subhdr{5.2 七步演算，每一步说清依据}

\textbf{第 1 步（换标定）}：把幅值换成有效值。由 §1.2 的 $I=I_m/\sqrt2$，

\[ I_1=\frac{12.7\sqrt2}{\sqrt2}=12.7\,\mathrm{A} ,\qquad
   I_2=\frac{11\sqrt2}{\sqrt2}=11\,\mathrm{A} . \]

\textbf{第 2 步（写相量）}：由 §3.5 的约定 $\dot I=I\angle\psi_i$，

\[ \dot I_1=12.7\angle30^\circ\,\mathrm{A} ,\qquad \dot I_2=11\angle-60^\circ\,\mathrm{A} . \]

\textbf{第 3 步（判定可加）}：两者的 $\omega$ 都是 $314\,\mathrm{rad/s}$，
满足 §3.3 的同频前提，故 $\dot I=\dot I_1+\dot I_2$ 合法。

\textbf{第 4 步（换成直角）}：加法只能用直角形式（§4.2）。逐个展开：

\[ \dot I_1=12.7\cos30^\circ+j\,12.7\sin30^\circ=11.0+j6.35 , \]
\[ \dot I_2=11\cos(-60^\circ)+j\,11\sin(-60^\circ)=5.5-j9.53 . \]

这里 p26 写的是「$11(\cos60^\circ-j\sin60^\circ)$」，用的是偶奇性
$\cos(-60^\circ)=\cos60^\circ$、$\sin(-60^\circ)=-\sin60^\circ$，数值一样。

\textbf{第 5 步（实部虚部分别相加）}：

\[ \dot I=(11.0+5.5)+j(6.35-9.53)=16.5-j3.18\,\mathrm{A} . \]

\textbf{第 6 步（换回极坐标）}：由 §4.1 的两条改写式，

\[ I=\sqrt{16.5^2+3.18^2}\approx16.8\,\mathrm{A} ,\qquad
   \psi_i=\arctan\frac{-3.18}{16.5}\approx-10.9^\circ , \]

即 $\dot I\approx16.8\angle-10.9^\circ\,\mathrm{A}$。注意实部为正，$\arctan$ 直接可用（§4.1 的警告在此处不触发）。

\textbf{第 7 步（还原瞬时式）}：由 §3.5，把 $\sqrt2$ 与 $\omega t$ 补回去：

\[ i=16.8\sqrt2\sin(314t-10.9^\circ)\,\mathrm{A} . \]

有效值 $I=16.8\,\mathrm{A}$——这正是 p26 右下角红框里那个数。

\aside{第 1 步用的是 §1.2 的 $I=I_m/\sqrt2$（跨小节：那在 p9，本例在 p26）；
第 4、6 步用的是 §4.1 的两条改写式（p20--p21）；
第 7 步用的是 §3.5 的还原约定（p25）。三处前提都不在本页，故在此一并点出。
代入本题：$12.7\sqrt2/\sqrt2=12.7$；$16.5-j3.18$ 的模是 $16.8$、角是 $-10.9^\circ$；
$16.8$ 是该乘 $\sqrt2$ 才成为幅值的有效值。}

\subhdr{5.3 独立校验：两相量恰差 $90^\circ$，构成直角三角形}

课件的 $16.8$ 与 $-10.9^\circ$ 不是我算出来就算完，最好有一条\textbf{不依赖课件中间数}的旁证。
看两个初相：$30^\circ-(-60^\circ)=90^\circ$。\textbf{两个相量恰好垂直}，于是合成量就是直角三角形的斜边：

\[ |\dot I|=\sqrt{12.7^2+11^2}=\sqrt{161.29+121}=\sqrt{282.29}\approx16.80\,\mathrm{A} . \]

角度也能独立算：以 $\dot I_2$ 的 $-60^\circ$ 为基准，$\dot I_1$ 在它前面 $90^\circ$，故

\[ \psi_i=-60^\circ+\arctan\frac{12.7}{11}=-60^\circ+49.10^\circ=-10.90^\circ . \]

两条都落在课件的 $16.8$ 与 $-10.9^\circ$ 上。\textbf{这条捷径只在相位差为 $90^\circ$ 时可用}，
但它给出的旁证与课件的直角坐标加法完全独立，是可靠的第二意见。

\subhdr{5.4 瞬时值层面的验证：相量相加 $\equiv$ 逐时刻相加}

§3.3 已经证明「同频时相量相加等价于瞬时值相加」。等价不是口头承诺，可以直接代两个时刻验：

\par\vspace{2pt}
\noindent\begin{tabular}{@{}p{1.9cm}p{3.0cm}p{3.0cm}p{3.4cm}p{3.7cm}@{}}
\tabtop
\thead{$t$} & \thead{$i_1$ / A} & \thead{$i_2$ / A} & \thead{$i_1+i_2$ / A} & \thead{由 $16.8\sqrt2\sin$ 算 / A}
\\
\tabmid
$0$ & $8.98$ & $-13.47$ & $-4.49$ & $-4.49$
\\
$0.005\,\mathrm{s}$ & $15.56$ & $7.77$ & $23.33$ & $23.32$
\tabbot
\end{tabular}

两列在三位有效数字内一致；第二位小数的差别来自 $16.8$ 与 $-10.9^\circ$ 的四舍五入。

\subhdr{5.5 为什么 $23.7$ 是错的：三层解释}

把两个有效值直接相加得 $12.7+11=23.7\,\mathrm{A}$，比真值 $16.8\,\mathrm{A}$ 大了四成。
必须说清它错在哪一层，否则下次还会犯。

\begin{itemize}
\item \textbf{几何层}：$12.7$ 与 $11$ 是两支\textbf{不同方向}的箭头的长度。长度相加只在两箭头同向时等于合成长度；
本例夹角 $90^\circ$，合成长度必然小于长度之和（三角形两边之和大于第三边）。
\item \textbf{代数层}：$|\dot I_1+\dot I_2|=|16.5-j3.18|=16.8$，而 $|\dot I_1|+|\dot I_2|=12.7+11=23.7$。
不等式 $|\dot I_1+\dot I_2|\le|\dot I_1|+|\dot I_2|$ 就是三角不等式，
它取等号的条件是两相量同向，即 $\varphi=0$。
\item \textbf{数值层}：若真按 $23.7$ 写瞬时式，初相就没有着落了——
$23.7\sqrt2\sin(314t+?)$ 里的问号填什么都不对，因为 $23.7$ 根本没有对应的相量。
\end{itemize}

\warnbox{\textbf{一句话判据}：\textbf{同相才能代数相加}。p30 的电阻上 $u$ 与 $i$ 同相，
所以 $U=RI$ 里 $U$ 与 $I$ 是「同向」的，才能直接乘一个 $R$；
p35 的电感与电容电压与总电压不同相，所以 $U\ne U_R+U_L+U_C$（§7.2 给正确式）。}

\subhdr{5.6 用相量图画一遍（p26 没有画，这里补）}

相量相加有两种作图法，结果相同：

\begin{itemize}
\item \textbf{平行四边形法}：两支相量共起点，作平行四边形，对角线即合成量。
\item \textbf{三角形法}：把 $\dot I_2$ 的起点平移到 $\dot I_1$ 的箭头端，从原点连到 $\dot I_2$ 的箭头端即合成量。
\end{itemize}

本例用三角形法最好看：$\dot I_1$ 从原点指向 $30^\circ$ 方向、长 $12.7$；
从它末端再画 $\dot I_2$，方向 $-60^\circ$、长 $11$。两段夹角 $90^\circ$，
所以合成量就是直角三角形的斜边 $16.8$，方向由「$-60^\circ$ 加 $49.1^\circ$」定出，落在 $-10.9^\circ$。
§5.3 的校验与这张图是同一件事的两种说法。

\aside{考场上的用法：先把两个相量按比例画出来，看合成量落在哪两个分量的方向之间。
本例 $-10.9^\circ$ 落在 $-60^\circ$ 与 $30^\circ$ 之间，这是必然的（分量的凸组合）——
一旦你算出的角不在这两个角之间，立刻回查第 4、5 步的符号。}

%=============================================================
%=============================================================
\hdr{§6}{单元件伏安关系：R、L、C 的相量式是怎么来的}

\subhdr{6.1 统一套路：三步，三个元件各走一遍}

p29 的提问是「If AC voltage is applied across a resistor, an inductor or a capacitor, what is the
relationship between the voltage and the current?」，Objectives 里补了一句
「especially the phase shift between the voltage and the current」。
p30--p32 三页各给一张波形图、一张相量图、两条式子，\textbf{但三页都没有推导}。
三个元件可以用完全相同的三步走完：

\begin{enumerate}
\item \textbf{设一个参考电流}：取 $i=\sqrt2I\sin\omega t$（即令 $\psi_i=0$）。
元件伏安关系是「绝对」的，与初相无关，所以把参考相取成 $\psi_i=0$ 不失一般性；
求出来的 $\varphi$ 就是该元件固有的相位差。
\item \textbf{用时域关系求 $u(t)$}：$R$ 用 $u=Ri$，$L$ 用 $u=L\dd i/\dd t$，$C$ 用 $i=C\dd u/\dd t$（电容先反解 $u$）。
\item \textbf{读两个数}：把 $u(t)$ 化成 $\sqrt2U\sin(\omega t+\varphi)$ 的标准形，
前面那个系数除以 $\sqrt2$ 就是 $U$，多出来的角就是 $\varphi$，相量即 $\dot U=U\angle\varphi$。
\end{enumerate}

\textbf{参考方向先定住}：p29--p32 三页画的是同一套约定——\textbf{电压 $u$ 上正下负，电流 $i$ 从正端流入}
（即 $u$ 与 $i$ 取关联参考方向）。下面所有 $+90^\circ$、$-90^\circ$ 的符号都以这套约定为准；
若某题把 $i$ 的箭头画反了，所有的相位差都要各加 $180^\circ$（超前与滞后互换）。
这不是细节，而是本题符号的唯一来源。

\subhdr{6.2 电阻：同相，$\dot U=R\dot I$（p30）}

第 1 步取 $i=\sqrt2I\sin\omega t$；第 2 步直接代入 $u=Ri$：

\[ u=\sqrt2\,(RI)\sin\omega t . \]

已经标准形，第 3 步读出 $U=RI$、$\varphi=0$——\textbf{电压与电流同相}。
相量写成

\[ \dot U=R\dot I ,\qquad U=RI . \]

\textbf{为什么同相}：电阻的关系 $u=Ri$ 是代数关系，不含时间导数，
电流取什么值，电压\textbf{同一瞬间}就是它的 $R$ 倍，所以两条曲线同时过零、同时到峰。
p30 右边那条绿 $u$ 与褐 $i$ 完全重合的波形图就是这个意思。

\warnbox{\textbf{p30 的提问没有答案}：「Is the reading of the Multimeter peak value or effective
value?」答案是\textbf{有效值}，理由见 §1.3。这条与 p9 的第二问是同一个问题，出现了两次，
说明它是老师反复想让人记住的点。}

\subhdr{6.3 电感：电压超前 $90^\circ$，$\dot U=jX_L\dot I$（p31）}

p31 全页\textbf{找不到} $u=L\dfrac{\dd i}{\dd t}$——它直接给了 $\dot U=jX_L\dot I$ 与
$X_L=\omega L=2\pi fL$，然后在右侧问了一句「Why voltage leads current by ninety degree?」，
\textbf{没有答案}。这一问一答要把整条链补齐。

第 2 步代入 $u=L\dfrac{\dd i}{\dd t}$，其中 $\dfrac{\dd i}{\dd t}=\sqrt2I\omega\cos\omega t$：

\[ u=L\cdot\sqrt2I\omega\cos\omega t=\sqrt2\,(\omega LI)\cos\omega t . \]

第 3 步把 $\cos$ 化成正弦（因为本册以 $\sin$ 为基准，见 §2.1 与 §3.1）：

\[ \cos\omega t=\sin\left(\omega t+90^\circ\right)
   \quad\Longrightarrow\quad
   u=\sqrt2\,(\omega LI)\sin\left(\omega t+90^\circ\right) . \]

读出两个数：

\[ U=\omega LI=X_LI ,\qquad \varphi=+90^\circ ,\qquad X_L=\omega L=2\pi fL . \]

$\varphi=+90^\circ$ 就是「\textbf{电压超前电流 $90^\circ$}」。写成相量：$\dot U=jX_L\dot I$——
那个 $j$ 正是 §4.4 的旋转因子，乘 $j$ 就是超前 $90^\circ$。

\aside{跨三页的一条链：p31 的 $\dot U=jX_L\dot I$ 里的 $j$，就是 p27 的 $e^{j90^\circ}=j$；
而 $90^\circ$ 这个数来自上式 $\cos\omega t=\sin(\omega t+90^\circ)$。
代入本题：$X_L=\omega L$，所以 $\dot U_L=j\omega L\dot I$，模 $U_L=\omega LI$，角 $\psi_u=\psi_i+90^\circ$。}

\textbf{回答 p31 那个提问（物理解释）}：电感电压正比于\textbf{电流的变化率}。
电流过零的那一刻变化最快，所以电压此刻最大；电流到峰值的那一刻变化率为零，所以电压此刻为零。
因此\textbf{电压的峰比电流的峰早四分之一周期出现}——早四分之一周期，就是超前 $90^\circ$。
不是「电感故意让电压早到」，而是「电压盯着变化率、变化率比电流本身早四分之一周期」。

\textbf{数值例}：$L=0.1\,\mathrm{H}$、$f=50\,\mathrm{Hz}$，则
$X_L=2\pi\times50\times0.1\approx31.4\,\Omega$；若 $I=2\,\mathrm{A}$，则
$U_L=X_LI\approx62.8\,\mathrm{V}$，且 $\psi_u=\psi_i+90^\circ$。

\subhdr{6.4 电容：电流超前 $90^\circ$，$\dot U=-jX_C\dot I$（p32）}

电容与电感互为镜像，但 p32 有两处硬错，必须先钉住再讲。

\textbf{第一处：问题抄错了。}p32 右上的提问仍是「Why voltage leads current by ninety degree?」，
这句话是 p31 电感页的。\textbf{电容上恰恰相反：电流超前电压 $90^\circ$}，
也就是电压\textbf{滞后}电流 $90^\circ$。同一页的波形图（红线 $i$ 的峰在左、绿线 $u$ 的峰在右）
与相量图（$\dot I$ 竖直向上、$\dot U$ 水平向右，$\dot I$ 比 $\dot U$ 靠前 $90^\circ$）都画对了，
只有这句问句是错的——它是从上一页直接复制过来的。

\textbf{第二处：相量式的下标印错了。}p32 印的是 $\dot U=-jX_L\dot I$，
而容抗记作 $X_C$、同页正文写的就是 $X_C=1/(\omega C)$。正确的式子只能是

\[ \dot U=-jX_C\dot I ,\qquad X_C=\frac{1}{\omega C}=\frac{1}{2\pi fC} . \]

这个笔误不是孤例：p33 的汇总表、p38 与 p39 的表格里电容那一格同样印成 $-jX_L\dot I$。
\textbf{四处同错，说明是模板复制造成的系统性笔误}，不是排版掉落。
本讲义全篇按 $-jX_C$ 走。

\textbf{推导。}取 $u=\sqrt2U\sin\omega t$（以电压为参考更方便，因为电容给的是 $i=C\dd u/\dd t$）：

\[ i=C\frac{\dd u}{\dd t}=C\cdot\sqrt2U\omega\cos\omega t
   =\sqrt2\,(\omega CU)\cos\omega t=\sqrt2\,(\omega CU)\sin\left(\omega t+90^\circ\right) . \]

读出两个数：

\[ I=\omega CU\ \Longrightarrow\ U=\frac{I}{\omega C}=X_CI ,\qquad
   \psi_i=\psi_u+90^\circ . \]

$\psi_i=\psi_u+90^\circ$ 即「\textbf{电流超前电压 $90^\circ$}」，等价于「电压滞后电流 $90^\circ$」。
相量形式：$\dot U=\dfrac{\dot I}{j\omega C}=-j\dfrac{\dot I}{\omega C}=-jX_C\dot I$。
乘 $-j$ 是顺时针转 $90^\circ$（§4.4），恰好对上「电压滞后」。

\aside{这里用了一次 $1/j=-j$：$\dfrac{1}{j\omega C}=\dfrac{1}{j}\cdot\dfrac{1}{\omega C}
=-jX_C$。验算：$j\cdot(-j)=1$，故 $-j=1/j$。这个负号不是规定出来的，
是 $1/j$ 算出来的——这也是 §4.4「乘 $-j$ 转 $-90^\circ$」与 §6.3「乘 $j$ 转 $+90^\circ$」的差别所在。}

\textbf{回答那个（被抄错的）提问的正确版本——「为什么电容上电流超前电压 $90^\circ$」}：
电容电流正比于\textbf{电压的变化率}。电压过零那一刻变化最快，所以电流此刻最大；
电压到峰值那一刻变化率为零，所以电流此刻为零。于是\textbf{电流的峰比电压的峰早四分之一周期}。
与电感完全对称：电感是「电压盯变化率」，电容是「电流盯变化率」。

\textbf{数值例}：$C=31.8\,\mu\mathrm{F}$、$f=50\,\mathrm{Hz}$，则
$X_C=1/(2\pi\times50\times31.8\times10^{-6})\approx100\,\Omega$；若 $U=220\,\mathrm{V}$，
则 $I=220/100=2.2\,\mathrm{A}$，且 $\psi_i=\psi_u+90^\circ$。

\subhdr{6.5 感抗与容抗的频率特性：一张表看两头}

把 §6.2--§6.4 的三条写成阻抗，顺便看 $\omega\to0$ 与 $\omega\to\infty$ 两个极限：

\par\vspace{2pt}
\noindent\begin{tabular}{@{}p{1.2cm}p{4.0cm}p{4.7cm}p{4.7cm}@{}}
\tabtop
\thead{元件} & \thead{相量式} & \thead{直流（$f\to0$）} & \thead{高频（$f\to\infty$）}
\\
\tabmid
$R$ & $\dot U=R\dot I$ & $U=RI$，不变 & 不变
\\
$L$ & $\dot U=j\omega L\dot I$ & $X_L\to0$：相当于\textbf{短路} & $X_L\to\infty$：相当于\textbf{开路}
\\
$C$ & $\dot U=-\dfrac{j}{\omega C}\dot I$ & $X_C\to\infty$：相当于\textbf{开路} & $X_C\to0$：相当于\textbf{短路}
\tabbot
\end{tabular}

这张表是后面几章（滤波、耦合电容、旁路电容、变压器）的伏笔，本册只立在这里不展开。
一个直观记法：\textbf{电感「通直流、阻交流」，电容「通交流、阻直流」}。

\subhdr{6.6 p33 那张汇总表：补上时域关系一列，并订正下标}

p33 用四行三列把三个元件并排。它的四行是 Waveform、Phasor diagram、Effective expression、Phasor
expression；\textbf{缺了最该有的一行：时域关系}。补上之后，整张表才回答得了「式子从哪来」：

\par\vspace{2pt}
\noindent\begin{tabular}{@{}p{3.3cm}p{5.3cm}p{3.4cm}p{2.6cm}@{}}
\tabtop
\thead{项目} & \thead{电阻 $R$} & \thead{电感 $L$} & \thead{电容 $C$}
\\
\tabmid
时域关系 & $u=Ri$ & $u=L\dfrac{\dd i}{\dd t}$ & $i=C\dfrac{\dd u}{\dd t}$
\\
$u$ 与 $i$ 的相位 & 同相（$\varphi=0$） & $u$ 超前 $i$ $90^\circ$ & $u$ 滞后 $i$ $90^\circ$
\\
有效值式 & $U=RI$ & $U=X_LI$ & $U=X_CI$
\\
$X$ 的表达式 & —— & $X_L=\omega L=2\pi fL$ & $X_C=\dfrac{1}{\omega C}=\dfrac{1}{2\pi fC}$
\\
相量式 & $\dot U=R\dot I$ & $\dot U=jX_L\dot I$ & $\dot U=-jX_C\dot I$
\tabbot
\end{tabular}

\aside{第三行第三列（p33 的电容相量式）原印 $--jX_L\dot I$，本表按 $X_C$ 订正，
理由见 §6.4。第三行第二列的「$u$ 滞后 $i$」与 p32 的问句冲突时，以波形图与相量图为准——
那两张图两路读法（文字层、读图记录）一致，而问句只有一处。}

%=============================================================
\hdr{§7}{相量形式的 KVL 与串联 RLC：$Z$ 是怎么长出来的}

\subhdr{7.1 相量 KVL 凭什么成立（p36 的第一行）}

p36 第一行写「KVL in phasor form：$\dot U=\dot U_R+\dot U_L+\dot U_C$」，\textbf{这一步没有解释}。
它成立靠三件事：瞬时 KVL 本来就成立、三项同频、以及 $\operatorname{Im}$ 的线性性。

时域里，串联电路共用同一个电流 $i$，KVL 对\textbf{每一个时刻}都成立：

\[ u(t)=u_R(t)+u_L(t)+u_C(t) . \]

三项都是与 $i$ 同频（$\omega$）的正弦量，可以按 §3.3 那套写法写成取虚部。
每一项各有一个相量，于是

\[ u(t)=\operatorname{Im}\left[\sqrt2\left(\dot U_R+\dot U_L+\dot U_C\right)e^{j\omega t}\right]
      =\operatorname{Im}\left[\sqrt2\,\dot U\,e^{j\omega t}\right] , \]

比较两边的相量部分即得 $\dot U=\dot U_R+\dot U_L+\dot U_C$。
\textbf{前提是两个}：三项同频（否则提不出公因子，§3.3）、电路是线性集中参数的
（否则 $u_R$ 等不再是同频正弦）。这两条全册默认成立，从未写出来——§10 记一条。

\subhdr{7.2 p35 那道题：$3\,\mathrm{V}$ 错在哪，正确答案是 $1\,\mathrm{V}$}

p35 的题面是「If $U_R=1\,\mathrm{V}$, $U_L=1\,\mathrm{V}$, $U_C=1\,\mathrm{V}$, then find the total
voltage」，旁边给出 $U=U_R+U_L+U_C=3\,\mathrm{V}$ 并打了一个红叉。\textbf{红叉后没有正确答案}，
答案要等到 p37 的相量图。把 p37 的图翻译成式子：

三个电压相量是 $\dot U_R=R\dot I$（与 $\dot I$ 同向）、$\dot U_L=jX_L\dot I$（超前 $90^\circ$）、
$\dot U_C=-jX_C\dot I$（滞后 $90^\circ$）。\textbf{$\dot U_L$ 与 $\dot U_C$ 方向正相反}，
当 $U_L=U_C=1\,\mathrm{V}$ 时两者等长反向，相加为零：

\[ \dot U_L+\dot U_C=jX_L\dot I-jX_C\dot I
   =j\left(X_L-X_C\right)\dot I=j(1-1)\dot I=0 . \]

于是

\[ \dot U=\dot U_R+\dot U_L+\dot U_C=\dot U_R ,\qquad U=U_R=1\,\mathrm{V} . \]

\textbf{正确答案是 $1\,\mathrm{V}$，不是 $3\,\mathrm{V}$。}三个有效值相等，合成却只等于其中一个——
这看起来反直觉，但相量图上一目了然：两个竖直的箭头等长反向，互相抵消，只剩水平那一个。

\warnbox{\textbf{一般式（电压三角形）}：三个电压有效值与总电压的关系是
\[ U=\sqrt{U_R^{2}+\left(U_L-U_C\right)^{2}} , \]
\textbf{不是} $U=U_R+U_L+U_C$。代本题：$\sqrt{1^2+(1-1)^2}=1\,\mathrm{V}$。
代 p37 的读法：$\sqrt{1^2+0}=1\,\mathrm{V}$。这个式子 §7.5 会给出严格来历。
\textbf{[教材补充]}：这条显式关系 p35--p37 没有印出来（p37 只有图），
本讲义由 p36 的相量 KVL 与 p42 的相量图合成，补它是为了让「$1\,\mathrm{V}$」有一个可代入的算式。}

\subhdr{7.3 从三项相加到 $Z$：p36 那三行标注在做什么}

p36 把推导压成三行，右边各贴一个标签。逐行展开：

\textbf{第一行（KVL in phasor form）}：$\dot U=\dot U_R+\dot U_L+\dot U_C$，§7.1 已证。

\textbf{第二行（Current expression）}：把 §6 的三条元件关系代进去。串联的要点是\textbf{共用同一个 $\dot I$}：

\[ \dot U=R\dot I+jX_L\dot I-jX_C\dot I . \]

\textbf{第三行（U-I Relationship in phasor form）}：三项都含 $\dot I$，提出来：

\[ \dot U=\left[R+j\left(X_L-X_C\right)\right]\dot I . \]

方括号里那个量，p36 用大括号括起来并标上「Impedance: $Z$」。这就是阻抗的出生地：
\textbf{它不是新定义，是从 KVL 里提出公因子提出来的}。

\aside{串联才有「共用同一个 $\dot I$」这一步——这是「串联」二字的全部含义。
并联时共用的是电压，「提出公因子」的对象要换成 $\dot U$，公式形状随之改变（C2-2 §4 讲）。
代入本题：$R=30\,\Omega$、$X_L=80\,\Omega$、$X_C=40\,\Omega$ 时方括号 $=30+j40\,\Omega$，见 §7.7。}

\subhdr{7.4 阻抗的两个写法（p41）}

p41 先把 $Z$ 写成直角形式，再一步跳到极坐标。中间那条改写式就是 §4.1 的两条公式：

\[ Z=R+jX=R+j\left(X_L-X_C\right) ,\qquad X=X_L-X_C . \]
\[ |Z|=\sqrt{R^{2}+X^{2}} ,\qquad \varphi=\arctan\frac{X}{R} ,\qquad Z=|Z|\angle\varphi . \]

读法三条：

\begin{itemize}
\item $|Z|$ 永远是\textbf{非负数}，单位 $\Omega$；负号只能藏在 $\varphi$ 里。
\item $\varphi$ 的符号由 $X$ 的符号定，$X>0$ 时 $\varphi>0$，$X<0$ 时 $\varphi<0$。
\item \textbf{$Z$ 不是相量，不加点}。相量代表一个正弦量，$Z$ 只是一个复数比（$\dot U/\dot I$），
它本身不对应任何随时间变化的量。p41 的方框里写的是 $Z$，没有点，就是这个道理。
\end{itemize}

\subhdr{7.5 $U=|Z|I$ 是怎么来的：两边取模（p40）}

p40 给出两个式子并各配一句注解：$\dot U=Z\dot I$ 是「ohm's law in phasor form」，
含「Two kinds of information, including magnitude and phase」；
$U=|Z|I$ 是「Effective expression」，只含「the information of magnitude」。
\textbf{从前者到后者只做了一件事：两边取模}。

\[ \dot U=Z\dot I
   \quad\Longrightarrow\quad
   |\dot U|=|Z|\,|\dot I|
   \quad\Longrightarrow\quad
   U=|Z|I . \]

取模会丢掉相位（$Z$ 的角 $\varphi$ 与 $\dot I$ 的角 $\psi_i$ 都不见了），所以 $U=|Z|I$
\textbf{只能算大小}。要相位、要画相量图、要判断超前滞后，必须回到 $\dot U=Z\dot I$。
p40 那两句注解说的就是这件事，只是没把「取模」这个动作写出来。

\subhdr{7.6 阻抗角判三种电路（p42），以及一个边角}

p42 的三幅相量图给三种情形，判据只有一条：\textbf{看 $X$ 的符号}，也就是看 $\dot U_L$ 与 $\dot U_C$ 谁高。

\par\vspace{2pt}
\noindent\begin{tabular}{@{}p{2.6cm}p{2.5cm}p{3.4cm}p{6.4cm}@{}}
\tabtop
\thead{$X$ 的符号} & \thead{$\varphi$} & \thead{电路性质} & \thead{相量图上}
\\
\tabmid
$X>0$（$X_L>X_C$） & $0<\varphi<90^\circ$ & 感性（Inductive） & $U_L$ 竖直向上，比 $U_C$ 长；$\dot U$ 斜向右上
\\
$X<0$（$X_L<X_C$） & $-90^\circ<\varphi<0^\circ$ & 容性（Capacitive） & $U_C$ 向下，比 $U_L$ 长；$\dot U$ 斜向右下
\\
$X=0$（$X_L=X_C$） & $\varphi=0^\circ$ & 阻性（Resistive） & $U_L$ 与 $U_C$ 等长反向抵消；$\dot U$ 与 $\dot I$ 重合
\tabbot
\end{tabular}

\warnbox{\textbf{边角：$R=0$ 时 $\arctan(X/R)$ 失效}。纯电感（$R=0$、$X>0$）与纯电容（$R=0$、$X<0$）
两种情形下 $X/R$ 无定义，$\varphi$ 必须单独写 $\varphi=+90^\circ$ 或 $\varphi=-90^\circ$。
p41、p43 都只给 $\varphi=\arctan(X/R)$，没写这个前提——§10 记一条（适用区间）。
实际电路中 $R$ 往往很小但非零，此时 $\varphi$ 会非常接近 $\pm90^\circ$，数值上要小心精度。}

\aside{[待确认] p42 三幅相量图的左右顺序：文字层抽出的顺序是「$\varphi=0^\circ$ / $-90^\circ<\varphi<0^\circ$ /
$0<\varphi<90^\circ$」，下方性质标注抽出的顺序是「Capacitive / Resistive / Inductive」，
两者不一致（分层文本的抽取顺序不可靠）。读图记录给的是左 Resistive、中 Inductive、右 Capacitive。
本讲义按\textbf{判据}讲，不依赖这三幅图的左右次序；若需按图对号，请以页面图为准。}

\subhdr{7.7 一个完整的数值算例（本讲义自设，$[\text{教材补充}]$）}

课件到 p43 为止\textbf{没有一道完整的串联 RLC 数值例题}（p35/p37 只有 $1\,\mathrm{V}$ 的定性题）。
为了把 §7.1--§7.6 串起来，这里自设一组数，全部用本册已有的公式。
\textbf{[教材补充]}之所以补：不代一个数，$Z$、$\varphi$、$U=|Z|I$ 三条都还停在符号上。
设 $R=30\,\Omega$、$X_L=80\,\Omega$、$X_C=40\,\Omega$、$U=220\,\mathrm{V}$、$f=50\,\mathrm{Hz}$。

\par\vspace{2pt}
\noindent\begin{tabular}{@{}p{2.2cm}p{4.4cm}p{8.4cm}@{}}
\tabtop
\thead{量} & \thead{值} & \thead{依据}
\\
\tabmid
$X$ & $80-40=40\,\Omega$ & $X=X_L-X_C$（p41）
\\
$Z$ & $30+j40\,\Omega$ & $Z=R+jX$（p36）
\\
$|Z|$ & $\sqrt{30^2+40^2}=50\,\Omega$ & $|Z|=\sqrt{R^2+X^2}$（p41）
\\
$\varphi$ & $\arctan(40/30)=53.13^\circ$ & $\varphi=\arctan(X/R)$（p41、p43）
\\
$I$ & $220/50=4.4\,\mathrm{A}$ & $U=|Z|I$（p40）
\\
$U_R$ & $30\times4.4=132\,\mathrm{V}$ & $U_R=RI$（p30）
\\
$U_L$ & $80\times4.4=352\,\mathrm{V}$ & $U_L=X_LI$（p31）
\\
$U_C$ & $40\times4.4=176\,\mathrm{V}$ & $U_C=X_CI$（p32）
\tabbot
\end{tabular}

\textbf{两处必须自己验的地方}：

（1）\textbf{三个有效值的代数和是 $132+352+176=660\,\mathrm{V}$，远大于电源 $220\,\mathrm{V}$}。
这不是算错，而是「有效值不能代数相加」的又一次现身。按电压三角形（§7.2）：

\[ U=\sqrt{U_R^{2}+\left(U_L-U_C\right)^{2}}
   =\sqrt{132^{2}+\left(352-176\right)^{2}}
   =\sqrt{132^{2}+176^{2}}=220\,\mathrm{V} , \]

正好回到电源电压。\textbf{这一验算同时证明了两件事}：式 $\sqrt{U_R^2+(U_L-U_C)^2}$ 是对的，
以及「代数相加」必错。

（2）\textbf{$U_L=352\,\mathrm{V}$ 比电源 $220\,\mathrm{V}$ 还高}。串联谐振电路里元件电压高于电源电压，
是正常现象（$U_L$ 与 $U_C$ 大部分互相抵消，只把差值留给电源）。这一点在做实验时会让初学者吓一跳。

\textbf{电路性质}：$X=40>0$，故 $\varphi=53.13^\circ>0$，属\textbf{感性}，电压超前电流 $53.13^\circ$。
若把电容换成 $X_C=100\,\Omega$，则 $X=-20\,\Omega$、$|Z|=\sqrt{900+400}\approx36.1\,\Omega$、
$\varphi=\arctan(-20/30)\approx-33.7^\circ$，变成\textbf{容性}——同一个电路，只换一个电容，性质就翻了个面。

\subhdr{7.8 阻抗三角形（p43）}

p43 的原句是「The total impedance of an RLC series circuit can be found using impedance triangle.
First of all, we have to consider the resistance and reactance of each component then put them into
the impedance triangle to find the total impedance of the RLC circuit.」
图是一个直角三角形：水平边 $R$、竖直边 $X$、斜边 $|Z|$，$\varphi$ 夹在 $R$ 与 $|Z|$ 之间。

它与 §7.2 的电压三角形是\textbf{相似三角形}：把电压三角形（$U_R$、$U_X$、$U$）每边除以 $I$，
就得到阻抗三角形（$R$、$X$、$|Z|$），角 $\varphi$ 不变。所以两个三角形的三条边成比例，
$\sqrt{}$ 与 $\arctan$ 的式子形状也一模一样：

\[ |Z|=\sqrt{R^{2}+X^{2}} ,\qquad \varphi=\arctan\frac{X}{R} ,\qquad
   \cos\varphi=\frac{R}{|Z|} . \]

\aside{最后那个 $\cos\varphi=R/|Z|$（功率因数）本册没有用到，此处只作为「三角形三条边的第三个关系」备案，
C2-2 讲功率时会用上。代入 §7.7 的数：$\cos53.13^\circ=30/50=0.6$，与 $R/|Z|$ 一致。}

\subhdr{7.9 特例建模：p38、p39 两张空表的填法}

p38 问「Can you relate the expression with the circuit on the right?」，
p39 的第二张表「Phasor expression」一行\textbf{整行空白}（读图记录确认未填）。
两张表问的是同一件事：\textbf{把 $R$、$X_L$、$X_C$ 中的若干个置零，一般式会退化成什么}。

\par\vspace{2pt}
\noindent\begin{tabular}{@{}p{3.4cm}p{4.6cm}p{6.9cm}@{}}
\tabtop
\thead{电路} & \thead{相量式} & \thead{从一般式怎么退化}
\\
\tabmid
只 $R$（p38 第 1 格） & $\dot U=R\dot I$ & 令 $X_L=X_C=0$
\\
只 $L$（p38 第 2 格） & $\dot U=jX_L\dot I$ & 令 $R=0$、$X_C=0$
\\
只 $C$（p38 第 3 格） & $\dot U=-jX_C\dot I$ & 令 $R=0$、$X_L=0$
\\
$R+L$（p39 第 1 格） & $\dot U=\left(R+jX_L\right)\dot I$ & 令 $X_C=0$
\\
$R+C$（p39 第 2 格） & $\dot U=\left(R-jX_C\right)\dot I$ & 令 $X_L=0$
\\
$L+C$（p39 第 3 格） & $\dot U=j\left(X_L-X_C\right)\dot I$ & 令 $R=0$
\tabbot
\end{tabular}

这就是 p38 那一问的答案：\textbf{不是「再记六个公式」，而是「一个公式置零六次」}。
一般式 $\dot U=[R+j(X_L-X_C)]\dot I$ 一旦记住，六种特例都是它的退化情形。

%=============================================================
\hdr{§8}{公式卡与解题五步}

\subhdr{8.1 公式卡（按做题顺序排，不按课件顺序）}

\par\vspace{2pt}
\noindent\begin{tabular}{@{}p{1.9cm}p{6.3cm}p{2.3cm}p{4.9cm}@{}}
\tabtop
\thead{量} & \thead{公式} & \thead{依据} & \thead{前提}
\\
\tabmid
$I$ & $I=I_m/\sqrt2$ & p9 & 正弦量
\\
$f,\ \omega$ & $\omega=2\pi f=2\pi/T$ & p7 & ——
\\
$\dot I$ & $\dot I=I\angle\psi_i$ & p26 & 有效值相量
\\
$i(t)$ & $i=\sqrt2I\sin(\omega t+\psi_i)$ & [C2-1 p25 反推] & 与上式配对，缺一不可
\\
$a+jb$ & $c=\sqrt{a^2+b^2}$，$\psi=\arctan(b/a)$ & p20（补） & $a>0$；$a<0$ 时加 $180^\circ$
\\
$\dot U_R$ & $R\dot I$ & p30 & 同上
\\
$\dot U_L$ & $jX_L\dot I$，$X_L=\omega L$ & p31 & 同上
\\
$\dot U_C$ & $-jX_C\dot I$，$X_C=1/(\omega C)$ & p32（订正） & 同上
\\
$Z$ & $R+j(X_L-X_C)$ & p36 & 串联
\\
$|Z|,\ \varphi$ & $\sqrt{R^2+X^2}$，$\arctan(X/R)$ & p41、p43 & $R>0$
\\
$U$ & $|Z|I$ & p40 & 只要大小
\\
$U$ & $\sqrt{U_R^2+(U_L-U_C)^2}$ & p42（补） & 串联、同频
\tabbot
\end{tabular}

\subhdr{8.2 五步流程（拿到任何一道交流题都照这个顺序）}

\begin{enumerate}
\item \textbf{读频率}：从 $314t$ 这类项里认出 $\omega$，换出 $f=314/(2\pi)\approx50\,\mathrm{Hz}$。
若题里两个正弦量 $\omega$ 不同，\textbf{停}——相量法不适用（§3.3）。
\item \textbf{定标定}：看 $\sin$ 前面有没有 $\sqrt2$，决定用有效值相量还是幅值相量，并在卷首写一句声明（§3.5）。
\item \textbf{换直角}：凡要加减，一律先化成 $a+jb$（§4.2）。
\item \textbf{用阻抗}：串联就写 $\dot U=Z\dot I$，然后视需要选极坐标（要相位）或取模（只要大小）。
\item \textbf{还原并自检}：把结果写回 $\sqrt2I\sin(\omega t+\psi)$；再做 §8.3 的三条自检。
\end{enumerate}

\subhdr{8.3 三条自检（交卷前 30 秒做完）}

\begin{itemize}
\item \textbf{量级自检}：$|Z|\ge R$ 且 $|Z|\ge|X|$；串联时 $U\ge U_R$；
若算出 $|Z|<R$，一定是 $X$ 的符号或平方算错了。
\item \textbf{相位自检}：$\varphi$ 的符号要与 $X$ 的符号一致（$X>0\Rightarrow\varphi>0$）；
合成的相量角必落在两个分量的角之间（§5.6）。
\item \textbf{$\sqrt2$ 自检}：最终写回瞬时式时，$\sin$ 前面\textbf{必须}有 $\sqrt2$。
若答案写成 $16.8\sin(\cdots)$，就漏了一个 $\sqrt2$（幅值应是 $16.8\sqrt2=23.76$）。
\end{itemize}

%=============================================================
\hdr{§9}{易错点（每条都指回某页某一步）}

\gapitem{（1）有效值与幅值混用}{p9 与 p25 之间隔着一条 $\sqrt2$ 却没有声明。
判据只有一条：\textbf{看 $\sin$ 前面有没有 $\sqrt2$}。p26 写 $12.7\angle30^\circ$ 对应
$12.7\sqrt2\sin(314t+30^\circ)$，模是有效值；若你写成 $12.7\sqrt2\angle30^\circ$，答案就差 $41\%$。}

\gapitem{（2）把 314 当成某个元件的参数}{p26 题面里的 $314$ 是 $\omega=2\pi\times50$，是频率，不是电阻也不是电感。
它出现在每个电流、每个电压的相位里，与元件无关。换成 $60\,\mathrm{Hz}$ 电网，全部换成 $377$。}

\gapitem{（3）有效值直接相加}{p35 的 $3\,\mathrm{V}$ 与 p26 的 $23.7\,\mathrm{A}$ 是同一个病。
正确做法：先换相量、再相加、最后取模。p35 的正确值是 $1\,\mathrm{V}$（§7.2）。}

\gapitem{（4）乘除时漏掉角的括号}{p22 印成 $c_1c_2\angle\psi_1+\psi_2$。正确是 $c_1c_2\angle(\psi_1+\psi_2)$。
另外除法是「角相减」且\textbf{顺序与被除数同}：$\psi_1-\psi_2$，不是 $\psi_2-\psi_1$。§4.3 有推导。}

\gapitem{（5）电容相量式抄成 $-jX_L\dot I$}{p32、p33、p38、p39 四处同错。容抗是 $X_C=1/(\omega C)$，
电感才叫 $X_L=\omega L$。看到电容的式子里出现 $X_L$，一律改 $X_C$（§6.4）。}

\gapitem{（6）超前与滞后记反}{p31 电感是 $u$ 超前 $i$（$\varphi=+90^\circ$），p32 电容是 $i$ 超前 $u$（$\varphi=-90^\circ$）。
两条互查办法：（a）波形图上\textbf{谁的峰在左边谁超前}；（b）相量图上\textbf{谁逆时针靠前谁超前}。
p32 的问句本身写反了，不能拿它当依据。}

\gapitem{（7）给 $Z$ 加点}{p36 的 $Z$ 是一个复数比，不是相量（p41 的方框里也没有点）。
加了点在概念上等于说「$Z$ 代表一个正弦量」，而它不代表。同理 $X_L$、$X_C$、$R$ 都不加点。}

\gapitem{（8）用 $\arctan(X/R)$ 不判情形}{p41、p43 只给 $\arctan(X/R)$。
两个坑：（a）$R=0$ 时无定义，必须单独写 $\varphi=\pm90^\circ$；
（b）若从 $a+jb$ 反求角而 $a<0$，要加 $180^\circ$（§4.1 的警告），否则会把第二象限的点算成第四象限。}

%=============================================================
\hdr{§10}{缺口清单：课件哪句话跳了步，本讲义补在哪}

\gapitem{机制缺口}{p9\quad 「Relationship between effective value and amplitude」，框里只有 $I=I_m/\sqrt2$。
\quad 缺：这个 $\sqrt2$ 从哪来（一个推导都没有）\quad 补：§1.2 热等效定义 + 两个积分}

\gapitem{答案形态}{p9\quad 「220V？Effective or amplitude (peak value)?」与
「Is the reading of the meters amplitude or effective value?」\quad 缺：两问都没有答案，p30 又问一次
\quad 补：§1.3（$220$ 与 $311$ 是同一电压；仪表读有效值）}

\gapitem{机制缺口}{p8\quad 「Why 50Hz/60Hz?」\quad 缺：整页只有这一问
\quad 补：§2.2 用工程折中与两套工业标准回答，并说明 314 的来历（标 [教材补充]）}

\gapitem{答案形态}{p2\quad 「Which one is not AC?」\quad 缺：只给「pulsing DC」一个标签，不指图号、不给判据
\quad 补：§1.0 按「是否跨横轴」逐图判定（并标 [待确认]）}

\gapitem{拼接缺口}{p27 与 p31/p32\quad 「$e^{j90^\circ}=j$」与「$\dot U=jX_L\dot I$」「$\dot U=-jX_C\dot I$」
\quad 缺：三页的 $j$ 是同一个旋转因子，没有任何一页把这条缝缝上
\quad 补：§0 第二处拼接、§4.4、§6.3 末尾}

\gapitem{适用区间}{p14\quad 「two sinusoidal AC current with the \textbf{same frequency}」
\quad 缺：为什么必须同频、不同频会怎样\quad 补：§3.3 取虚部的三行证明 + 拍频反例}

\gapitem{形式改写}{p25\quad 「Two kinds of information」（Magnitude / Phase），相量写成 $1+j=\sqrt2\angle45^\circ$
\quad 缺：相量与瞬时值的换算 $\dot I=I\angle\psi\iff i=\sqrt2I\sin(\omega t+\psi)$ 全册未写；
模取有效值还是幅值从未声明\quad 补：§3.5（约定 + 用 p26 反验 + $\sqrt2$ 判据）}

\gapitem{代数跳步}{p20\quad 「$\overline{Op}=a+jb=c(\cos\psi+j\sin\psi)=ce^{j\psi}=c\angle\psi$」
\quad 缺：$c=\sqrt{a^2+b^2}$ 与 $\psi=\arctan(b/a)$ 这两次改写一个都没写，
p21 的 $1+j\leftrightarrow\sqrt2\angle45^\circ$ 直接用它们\quad 补：§4.1（含 $a<0$ 补 $180^\circ$ 的警告）}

\gapitem{机制缺口}{p21\quad 「Why do we need so many forms?」\quad 缺：没有答案
\quad 补：§4.2 分工表，落成八字口诀「加减用直角，乘除用极坐标」}

\gapitem{形式改写}{p22\quad 「Multiplication: $A_1A_2=c_1c_2\angle\psi_1+\psi_2$」「Division：
$A_1/A_2=c_1/c_2\angle\psi_1-\psi_2$」\quad 缺：角的括号掉了，按字面读量纲都不合；
且两条法则没有推导\quad 补：§4.3（加括号 + 用指数形式相乘推出）}

\gapitem{机制缺口}{p27\quad 「Example: rotating factor $e^{j90^\circ}=j$」\quad
缺：为什么乘 $j$ 等于转 $90^\circ$；且页面上「Question:」框是空的、第二行右格没画箭头
\quad 补：§4.4（旋转因子的乘法法则解释 + 空题的补法，[待确认]）}

\gapitem{机制缺口}{p31\quad 「Why voltage leads current by ninety degree?」
\quad 缺：这一问没有答案；更关键的是\textbf{全页没有 $u=L\dd i/\dd t$}，$\dot U=jX_L\dot I$ 凭空出现
\quad 补：§6.3 三步推导（求导得 $\cos$、化 $\sin$、读出 $+90^\circ$）+ 物理说明}

\gapitem{形式改写 / 笔误}{p32\quad 「Why voltage leads current by ninety degree?」与「$\dot U=-jX_L\dot I$」
\quad 缺：电容是\textbf{电流}超前电压 $90^\circ$，问句从 p31 抄来写反了；相量式下标应为 $X_C$
（p33、p38、p39 同错，四处一致）\quad 补：§6.4 点名订正并重推；§6.6 给订正版汇总表}

\gapitem{适用前提}{p36\quad 「KVL in phasor form：$\dot U=\dot U_R+\dot U_L+\dot U_C$」
\quad 缺：瞬时 KVL 凭什么能直接变相量（同频 + 线性两条前提一字未写）
\quad 补：§7.1 用 $\operatorname{Im}$ 的线性性与「提出公因子」证明}

\gapitem{答案形态}{p35\quad 「$U=U_R+U_L+U_C=3\,\mathrm{V}$」被打红叉
\quad 缺：错了但没有正确答案，正确值 $1\,\mathrm{V}$ 只在 p37 靠图暗示
\quad 补：§7.2 相量相加得 $\dot U=\dot U_R$，并给一般式 $U=\sqrt{U_R^2+(U_L-U_C)^2}$}

\gapitem{形式改写}{p40\quad 「$U=|Z|I$ … Only the information of magnitude」
\quad 缺：这个式子是从 $\dot U=Z\dot I$ \textbf{两边取模}来的，取模这个动作没写，
读者会以为它是另一个独立结论\quad 补：§7.5（补出取模一行，并说明为什么丢了相位）}

\gapitem{适用区间}{p41、p43\quad 「$\varphi=\arctan(X/R)$」
\quad 缺：$R=0$（纯电感或纯电容）时无定义，必须单独写 $\pm90^\circ$；也没说象限
\quad 补：§7.6 的边角提示、§4.1 的象限警告}

\gapitem{答案形态}{p39\quad 第二张表「Phasor expression」一行整行空白
\quad 缺：$R+L$、$R+C$、$L+C$ 三种组合的相量式没有填
\quad 补：§7.9 给六格完整的填法（含 p38 的「一个公式置零六次」）}

\aside{§10 共 18 条，类型分布：机制缺口 5 条、答案形态 4 条、形式改写 4 条、适用区间/前提 3 条、
拼接缺口 1 条、代数跳步 1 条（其中 p32 一条兼属「形式改写」与「笔误」）。
按 §0 的观察，\textbf{「给了结果不给来历」是本册的系统性特征}（$\sqrt2$、$Z$、$U=|Z|I$、旋转因子都这样），
所以 §1、§6、§7 三节都按「先补推导、再回指页码」的方式写。}

%=============================================================
\hdr{§11}{三句话总结}

\begin{enumerate}
\item \textbf{有效值是热等效算出来的均方根}，等于幅值除以 $\sqrt2$；
\quad 一切仪表读数、铭牌、相量的模都是有效值，只有波形图顶端与 $\sin$ 前的系数才是幅值。
\item \textbf{相量能代替正弦量，靠的是「同频」与「取虚部的线性」}：
\quad 提出公因子 $e^{j\omega t}$ 之后只剩一个复数，它丢掉频率、留下模与初相；
\quad 加减用直角形式，乘除用极坐标形式，同相才能代数相加。
\item \textbf{元件伏安关系靠 $\dd/\dd t\leftrightarrow j\omega$ 变成阻抗}：
\quad 三个元件化成 $R$、$jX_L$、$-jX_C$，串联时 KVL 提出公因子得 $Z=R+j(X_L-X_C)$；
\quad 电压必须按相量相加，$U=\sqrt{U_R^2+(U_L-U_C)^2}$，永远不等于三个有效值之和。
\end{enumerate}

%=============================================================
\hdr{§12}{与课件的页码对照表}

\par\vspace{2pt}
\noindent\begin{tabular}{@{}p{1.6cm}p{7.4cm}p{6.6cm}@{}}
\tabtop
\thead{课件页} & \thead{课件内容} & \thead{本讲义位置}
\\
\tabmid
p1 & 封面：Chapter Two AC circuits & §0（封面不计内容）
\\
p2 & Definition of AC、五个波形、$question$「Which one is not AC?」 & §1.0
\\
p3--p4 & AC generator：磁铁、绕组、切割磁力线、碳刷滑环（p4 为放大图） & §1.0（发电机链条与正弦的来历）
\\
p5 & 直流单向 / 交流周期性换向 & §1.0
\\
p6 & $i=I_m\sin(\omega t+\psi)$、三要素、波形 & §2.1
\\
p7 & 周期、频率、角频率；$T=0.5\,\mathrm{s}$、$f=2\,\mathrm{Hz}$、$\omega=4\pi$ & §2.2
\\
p8 & $question$「Why 50Hz/60Hz?」 & §2.2
\\
p9 & 瞬时值/幅值/有效值、$I=I_m/\sqrt2$、两问、万用表照片 & §1.1--§1.3
\\
p10 & 相位与初相定义、E1（$i=10\sin(1000t+30^\circ)$） & §2.3
\\
p11 & E2（$u=311\sin(314t-60^\circ)$） & §2.3、§1.3
\\
p12 & E3：相位差 $\varphi=\psi_u-\psi_i=-90^\circ$ & §2.3
\\
p13 & 封面：2.2 Phasor representation & §0（封面不计内容）
\\
p14 & 提问（两个同频电流如何相加）、四个 Objective & §3.3
\\
p15 & 正弦波 $\leftrightarrow$ 旋转矢量（取纵轴投影） & §3.1
\\
p16 & 三种初相的波形与旋转矢量对照 & §3.1
\\
p17 & 旋转矢量 $\leftrightarrow$ 静止矢量 & §3.2
\\
p18 & 静止矢量 $\leftrightarrow$ 相量（含定义句） & §3.2
\\
p19 & 相量 $\leftrightarrow$ 复数 & §3.2、§4
\\
p20 & 复数四形式（直角、三角、指数、极坐标） & §4.1
\\
p21 & $1+j\leftrightarrow\sqrt2\angle45^\circ$ 对照表、$question$「Why so many forms?」 & §4.2
\\
p22 & 复数加减乘除四条 & §4.3
\\
p23 & E1、E2、E3 三道复数运算题（无答案） & §4.5
\tabbot
\end{tabular}

\par\vspace{2pt}
\noindent\begin{tabular}{@{}p{1.6cm}p{7.4cm}p{6.6cm}@{}}
\tabtop
\thead{课件页} & \thead{课件内容} & \thead{本讲义位置}
\\
\tabmid
p24 & 封面：2.2.2 Phasor representation of sinusoidal AC & §0（封面不计内容）
\\
p25 & $\dot I=1+j=\sqrt2\angle45^\circ$、$two kinds of information$ & §3.5
\\
p26 & Example3：$12.7\angle30^\circ+11\angle-60^\circ\Rightarrow16.8\angle-10.9^\circ$ & §5（主例题）
\\
p27 & 旋转因子 $e^{j90^\circ}=j$、$e^{-j90^\circ}=-j$、$question$（空） & §4.4
\\
p28 & 封面：2.3 Single parameter AC circuit & §0（封面不计内容）
\\
p29 & 提问（R/L/C 上电压与电流的关系）、$Objective$ & §6.1
\\
p30 & 纯电阻：波形同相、相量图、$U=RI$、$\dot U=R\dot I$、万用表提问 & §6.2、§1.3
\\
p31 & 纯电感：$u$ 超前 $90^\circ$、$X_L=\omega L$、$\dot U=jX_L\dot I$、$question$ & §6.3
\\
p32 & 纯电容：$i$ 超前 $90^\circ$、$X_C=1/(\omega C)$、$\dot U=-jX_L\dot I$（下标误） & §6.4（订正）
\\
p33 & 2.3.4 汇总表（波形、相量图、有效值式、相量式） & §6.6（补时域列）
\tabbot
\end{tabular}

\par\vspace{2pt}
\noindent\begin{tabular}{@{}p{1.6cm}p{7.4cm}p{6.6cm}@{}}
\tabtop
\thead{课件页} & \thead{课件内容} & \thead{本讲义位置}
\\
\tabmid
p34 & 封面：2.4 Series AC circuits & §0（封面不计内容）
\\
p35 & 提问 $U_R=U_L=U_C=1\,\mathrm{V}$、$3\,\mathrm{V}$ 打红叉 & §7.2
\\
p36 & 相量 KVL、$Current\ expression$、$Z$ 的出生地 & §7.1、§7.3
\\
p37 & 相量图作答：$U_L$ 与 $U_C$ 相消，$U=U_R$ & §7.2
\\
p38 & （1）特例建模、$question$「Can you relate\ldots」、三元件表 & §7.9
\\
p39 & （1）特例建模第二张表（$R+L$、$R+C$、$L+C$，相量式行空白） & §7.9
\\
p40 & （2）$\dot U=Z\dot I$ 与 $U=|Z|I$（两种信息） & §7.5
\\
p41 & （3）阻抗：直角式、极坐标式、$|Z|$ 与 $\varphi$ & §7.4
\\
p42 & （4）阻抗角：三种电路性质、$\varphi=\psi_u-\psi_i$ & §7.6（含 [待确认]）
\\
p43 & （5）阻抗三角形、$\varphi=\arctan(X/R)$ & §7.8
\tabbot
\end{tabular}

\flowerrule

{\small\color{inkgray}
本讲义所有数值已用 Python 独立复算：$220\sqrt2=311.1$、$50\,\mathrm{Hz}$ 对应 $\omega=314.16$、
$16.5-j3.18$ 的模为 $16.8015$、角为 $-10.897^\circ$、
$12.7^2+11^2$ 的平方根为 $16.80$、$-60^\circ+\arctan(12.7/11)=-10.90^\circ$、
E1 $=6+j7=9.22\angle49.40^\circ$、E2 $=\sqrt2\angle15^\circ=1.366+j0.366$、
E3 $=\sqrt2\angle75^\circ=0.366+j1.366$、
$X_L(0.1\,\mathrm{H},50\,\mathrm{Hz})=31.4\,\Omega$、$X_C(31.8\,\mu\mathrm{F},50\,\mathrm{Hz})=100.1\,\Omega$、
$|30+j40|=50$、$\arctan(40/30)=53.13^\circ$、$\sqrt{132^2+176^2}=220$。
标 [教材补充] 的三处（50/60 Hz 的由来、$U=\sqrt{U_R^2+(U_L-U_C)^2}$ 的一般式、§7.7 的完整算例）
不是课件原文，只作理解用，不作文档口径。}

\end{document}


